\documentclass[11pt]{article}
\usepackage[margin=1.1in]{geometry}
\usepackage{amsmath,amssymb,amsthm}
\usepackage{xcolor}
\usepackage{tocloft}
\usepackage{enumitem}
\setlist[enumerate]{label=(\arabic*),itemsep=2pt}
\setlength{\cftbeforesecskip}{2pt}
\usepackage[colorlinks=true,linkcolor=magenta,citecolor=magenta,urlcolor=magenta]{hyperref}

\theoremstyle{plain}
\newtheorem{theorem}{Theorem}[section]
\newtheorem{proposition}[theorem]{Proposition}
\newtheorem{lemma}[theorem]{Lemma}
\newtheorem{definition}[theorem]{Definition}
\newtheorem*{theoremA}{Theorem A}
\newtheorem*{theoremB}{Theorem B}
\newtheorem*{theoremC}{Theorem C}
\newtheorem*{theoremD}{Theorem D}
\newtheorem*{theoremE}{Theorem E}

\title{Italian projective geometry goes free}
\author{Claude O and Teo B}
\date{}

\begin{document}
\maketitle

\begin{abstract}
We develop basic free projective geometry, in the old style, with all the objects being given by explicit generators and relations. We construct free analogues of the Segre and Veronese embeddings, of the Grassmannians, flag manifolds and Stiefel manifolds, and of the determinantal varieties. The torus fixed points of the free full flag manifold are given by the quantum permutation group $S_N^+$, and the free Grassmannians are multiplicity free, with the same spectrum as the free projective spaces, and, at $N\geq3$, are not quotient spaces of $O_N^+$. We also discuss integration, with free Poisson laws appearing, the heat equation, with a threshold for the heat trace coming from the exponential growth of the multiplicities, the arithmetic aspects, and the relation with questions from particle physics.
\end{abstract}

\tableofcontents

\section*{Introduction}
\addcontentsline{toc}{section}{Introduction}

Classical projective geometry, as developed by the Italian school of Segre, Veronese, Castelnuovo, Enriques and Severi, is about explicit objects: projective spaces, products and the Segre embedding, Veronese embeddings and hypersurfaces, Grassmannians, flags, determinantal varieties. All these are given by concrete equations, and we refer to Harris \cite{har} and Shafarevich \cite{sha} for the theory. On the differential geometry side, the same objects appear as homogeneous spaces, with interesting harmonic analysis, see Helgason \cite{hel}.

\bigskip

In the free setting, the projective space $P^{N-1}_+$ is the abstract spectrum of the universal $C^*$-algebra generated by the entries of a projection $p=(p_{ij})_{i,j=1,\ldots,N}$ of trace 1, and its basic properties are discussed in the book \cite{ba1}. At the end of that discussion, the question of constructing the free Grassmannians, flag manifolds and Stiefel manifolds was raised. We discuss here this question, and more generally the free analogues of all the above classical objects.

\bigskip

Our viewpoint will be that of old-style algebraic geometry, with everything being given by simple generators and relations, and with the only modern ingredient being the freeness. The basic principle is that the free objects should be written in terms of projections, rather than in terms of determinants, which are not available in the free setting. Our first result, regarding the basic projective geometry constructions, is as follows:

\begin{theoremA}
We have free analogues of the Segre and Veronese embeddings, given by
$$P_{(ia),(jb)}=p_{ij}q_{ab}\quad,\quad
P_{IJ}=p_{i_1i_2}p_{i_3i_4}\ldots p_{j_2j_1}$$
with the Veronese coordinates being the products over the consecutive pairs of letters of the words $i_1\ldots i_dj_d\ldots j_1$. Also, the free Grassmannians, given by the projections of trace $L$, embed for $L\geq1$ into $P^{N^2-1}_+$, via $P_{(i_1i_2),(j_1j_2)}=\frac{1}{L}p_{i_1i_2}p_{j_2j_1}$.
\end{theoremA}

The last embedding is a substitute for the Pl\"ucker embedding, which is not available in the free setting, due to the lack of a free determinant. The Veronese embeddings allow, as classically, to view the hypersurfaces as hyperplane sections. Regarding symmetries, we prove that the projective quantum isometry groups of the real, complex and free Grassmannians are $PO_N\subset PU_N\subset PO_N^+=PU_N^+$, the last isomorphism being from \cite{bfu}, and the complex case being computer-assisted, via a Tannakian generation argument.

\bigskip

Our second result concerns the free flag manifolds $Fl^+_d$, with $d=(d_1,\ldots,d_k)$ being a composition of $N$, which appear from families of projections $q^1,\ldots,q^k\in M_N(A)$ summing up to 1, with $Tr(q^s)=d_s$. Equivalently, these are the free isospectral manifolds, of self-adjoint matrices $x=\sum_s\lambda_sq^s$ with prescribed spectrum. We have:

\begin{theoremB}
The free flag manifolds appear from the free orthogonal quantum group $O_N^+$, via
$$q^s_{ij}\to\sum_{r\in B_s}u_{ri}u_{rj}$$
with $B_1,\ldots,B_k$ being the blocks of sizes $d_1,\ldots,d_k$, and the real and complex orbit models of these spaces coincide. Moreover, the torus fixed points of the free full flag manifold are given by the quantum permutation group $S_N^+$.
\end{theoremB}

Here the fact that the real and complex versions coincide is a generalization of the isomorphism $PO_N^+=PU_N^+$, telling us that free projective geometry is ``scalarless''. As for the last assertion, this tells us that $S_N^+$ appears as a free Weyl group.

\bigskip

We will discuss as well the free determinantal varieties, via an incidence model, and the free Stiefel manifolds, which are the spaces of free partial isometries from \cite{ba1}, with $L$ rows. Our third result, which is the main one, concerns the harmonic analysis on the free Grassmannians. Recall that the irreducible corepresentations of $O_N^+$, with $N\geq2$, are $r_0=1$, $r_1=u$, and then $r_k\subset u^{\otimes k}$, given by the Jones-Wenzl projections, with $SU_2$ type fusion rules $r_k\otimes r_l=r_{|k-l|}+r_{|k-l|+2}+\ldots+r_{k+l}$, and with $\dim r_k=U_k(N)$, the Chebyshev polynomials of the second kind, see \cite{ba1}. Given a quantum space $X$ on which $O_N^+$ acts, its spectrum is the set of irreducible corepresentations appearing in $C(X)$, in the harmonic analysis sense, as for the spaces $L^2(G/H)$ in the classical case, and $X$ is called multiplicity free when each of them appears at most once. Classically, $Gr_{LN}$ is a symmetric space of rank $\min(L,N-L)$, and its spectrum is multiplicity free. In the free case, we have:

\begin{theoremC}
The free Grassmannian $Gr^+_{LN}$ with $1\leq L\leq N-1$, viewed inside $O_N^+$, is multiplicity free, with spectrum $r_0,r_2,r_4,\ldots$ independent of $L$, and so equal to the spectrum of $P^{N-1}_+$. Moreover, at $N\geq3$ it is not a quotient space of type $O_N^+/H$.
\end{theoremC}

Thus, the free Grassmannians behave like rank one spaces, and at $N\geq3$ are ``infinite dimensional'' in the Weyl law sense, as the free spheres, and we will see later that this produces a threshold for the heat trace, analogous to the Hagedorn temperature. For learning basic thermodynamics, Fermi \cite{fth} is what you need. Finally, we discuss integration, with the Weingarten formula for $O_N^+$ as main tool. A basic result here is as follows:

\begin{theoremD}
The coordinate $p_{11}\in C(Gr^+_{LN})$ has first moments $L/N$ and $L(L+1)/N(N+1)$, and with $N\to\infty$, $L$ being fixed, the variable $Np_{11}$ becomes free Poisson of parameter $L$.
\end{theoremD}

This is to be compared with the classical case, where $p_{11}$ follows the beta law $\nu_{L/2,(N-L)/2}$ in the real case, and $\nu_{L,N-L}$ in the complex case, and where with $N\to\infty$ the variable $Np_{11}$ becomes the chi-squared law $\chi^2_{L1}=\mu_{L/2,2}$, respectively the gamma law $\mu_{L1}=\chi^2_{2L,1/2}$. Here we use the following notations, for the beta and gamma laws:
$$\nu_{ab}=\frac{x^{a-1}(1-x)^{b-1}}{B(a,b)}\,dx\quad,\quad
\mu_{ms}=\frac{x^{m-1}e^{-x/s}}{s^m\Gamma(m)}\,dx$$
and $\chi^2_{nt}=\mu_{n/2,2t}$ is the law of $\sum_if_i^2$, with $f_1,\ldots,f_n$ being independent centered normal variables, of variance $t$.

\bigskip

In the second part of the paper we discuss the heat equation, the arithmetic aspects, and the relation with particle physics. Here the free projective spaces appear as free spaces of pure states, whose quantum isometries are the weak and color quantum symmetries from \cite{bop}, and the transition probabilities between pure states, which enter into Fermi's golden rule, have free analogues, as follows:

\begin{theoremE}
Given $2$ independent free full flags $q,r\in Fl^+_N$, with the integration coming from $O_N^+$, the free transition matrix $\theta_{st}=\sum_{ij}q^s_{ij}r^t_{ij}$ is bistochastic, with positive entries, and has the same law as the matrix $(u_{st}^2)_{st}$, where $u$ is the fundamental corepresentation of $O_N^+$.
\end{theoremE}

Classically, the relative position of 2 random real flags is given by a random orthogonal matrix $U$, and the transition matrix is the orthostochastic matrix $(U_{st}^2)$. In the free case the classical coset description fails, as explained above, but the same conclusion holds, with $O_N^+$ replacing $O_N$, and with the order of the indices in the definition of $\theta$ being crucial. We will see as well that the free mixing has an intrinsic CP-odd component, in the sense that a free analogue of the Jarlskog invariant has positive variance, in any dimension $N\geq3$. We also show that the free projective line is classical, $P^1_+=P^1_\mathbb C$, so that free projective geometry really starts at $N=3$.

\bigskip

Let us mention that noncommutative projective spaces, Grassmannians and flag manifolds have been studied from several other viewpoints, namely the $q$-deformed ones \cite{ddl}, \cite{dst}, \cite{fio}, \cite{hsz}, the toric deformations of Cirio, Landi and Szabo \cite{cls}, with noncommutative Pl\"ucker relations, the graded noncommutative projective spaces, with their Segre and Veronese morphisms \cite{ags}, \cite{bch}, and the quasi-Pl\"ucker coordinates \cite{lau}. In these approaches the noncommutativity comes from a deformation parameter, or from the use of a general division ring, as for the quasi-Pl\"ucker coordinates, while here it comes from freeness, with no parameter involved, and with the real and complex cases coinciding.

\bigskip

The paper is organized as follows: sections 1-3 deal with projective spaces, products, and Veronese embeddings, sections 4-6 deal with Grassmannians, flag manifolds, and determinantal varieties, and section 7 deals with Stiefel manifolds, spectral theory and integration. In sections 8-9 we discuss the heat equation and the Laplacians, following \cite{mfs}, \cite{dfw}, and the free real algebraic geometry and arithmetic aspects, following \cite{fba}, \cite{oba}. Finally, section 10, which is dedicated to Enrico Fermi, discusses the relation with particle physics, following \cite{bop}.

\section{Projective spaces}

We use the compact quantum space formalism from \cite{ba1}, and Woronowicz's compact quantum group formalism from \cite{wo1}. Our scalar products are linear at left, $<x,y>=\sum_ix_i\bar{y}_i$. All our matrices have entries in $C^*$-algebras, and $Tr$ denotes the usual trace, $Tr(p)=\sum_ip_{ii}$. We will need the following elementary fact, which will be used many times in what follows:

\begin{lemma}
\label{lem:trace}
If $q\in M_n(A)$ is a projection with $Tr(q)=0$, then $q=0$.
\end{lemma}

\begin{proof}
From $q=qq^*$ we get $q_{rr}=\sum_sq_{rs}q_{rs}^*\geq0$. Since these elements sum up to 0, they all vanish, and so $q_{rs}q_{rs}^*=0$ for any $r,s$, which gives $q=0$.
\end{proof}

Following \cite{ba1}, the free projective space is constructed as follows:

\begin{definition}
The free projective space $P^{N-1}_+$ is the abstract spectrum of the algebra
$$C(P^{N-1}_+)=C^*\left((p_{ij})_{i,j=1,\ldots,N}\Big|p=p^*=p^2,Tr(p)=1\right)$$
whose generators $p_{ij}$ are called projective coordinates.
\end{definition}

This algebra is well-defined, because the entries of a projection have norm $\leq1$. The commutative quotient is $C(P^{N-1}_\mathbb C)$, the characters corresponding to the rank one projections in $M_N(\mathbb C)$, and the real projective space $P^{N-1}_\mathbb R$ appears by further imposing $p=\bar{p}$.

\bigskip

The relation with the free spheres from \cite{ba1} is as follows. Given a point of the free real sphere, meaning self-adjoint variables $z_1,\ldots,z_N$ with $\sum_iz_i^2=1$, the matrix $p_{ij}=z_iz_j$ is a projection of trace 1, and similarly in the complex case, with $p_{ij}=z_i^*z_j$. We have:

\begin{proposition}
We have quotient maps $S^{N-1}_{\mathbb R,+}\to P^{N-1}_+$ and $S^{N-1}_{\mathbb C,+}\to P^{N-1}_+$, given by
$$p_{ij}=z_iz_j\quad,\quad p_{ij}=z_i^*z_j$$
and these are quotient maps, in the sense that the corresponding morphisms of algebras are injective.
\end{proposition}

\begin{proof}
The fact that we have projections of trace 1 follows from $\sum_iz_i^2=1$, and from $\sum_iz_iz_i^*=\sum_iz_i^*z_i=1$. For the injectivity, represent $C(P^{N-1}_+)$ faithfully on a Hilbert space $H$, set $K=pH^N$, and let $E_i:H\to H^N$ be the coordinate inclusions. With $T_i=pE_i:H\to K$ we have $T_i^*T_j=p_{ij}$, $\sum_iT_i^*T_i=1_H$ and $\sum_iT_iT_i^*=1_K$. Thus $z_i=(^{\ 0\ }_{T_i}{\ }^{T_i^*}_{\ 0\ })$ on $H\oplus K$ satisfy $z_i=z_i^*$ and $\sum_iz_i^2=1$, and the upper corner of $z_iz_j$ is $p_{ij}$. This corner is multiplicative on the algebra generated by the products $z_iz_j$, and recovers our faithful representation, so the real map is injective. In the complex case, the composition with the quotient map $S^{N-1}_{\mathbb C,+}\to S^{N-1}_{\mathbb R,+}$ is the real map, which gives the result.
\end{proof}

As a basic example, at $N=2$ the projective coordinates $p=(^a_{c^*}{\ }^c_b)$ of $P^1_+$, with $a=a^*$, $b=b^*$, $a+b=1$, satisfy the following relations, coming from $p^2=p$, the off-diagonal relation $ac+cb=c$ being equivalent to $ac=ca$:
$$cc^*=c^*c=ab\quad,\quad ac=ca$$

Thus $a$ is central and $c$ is normal, so the algebra $C(P^1_+)$ is commutative, and we have $P^1_+=P^1_\mathbb C=S^2$. In other words, free projective geometry really starts at $N=3$, in agreement with the fact that $PO_2^+=SO_3$ is classical. At $N\geq3$ the algebra $C(P^{N-1}_+)$ is no longer commutative, and is not even nuclear. Indeed, with $g_1,\ldots,g_N$ being the standard generators of $\mathbb Z_2^{*N}$, the matrix $p_{ij}=g_ig_j/N$ is a projection of trace 1, and this gives a quotient map onto the full group algebra of the even subgroup of $\mathbb Z_2^{*N}$, which is free of rank $N-1$.

In order to discuss symmetries, given a closed subgroup $G\subset O_N^+$, we say that $G$ acts on a closed subspace $X\subset P^{N-1}_+$ when the following formula defines a morphism of algebras:
$$\Phi:C(X)\to C(X)\otimes C(G)\quad,\quad
\Phi(p_{ij})=\sum_{ab}p_{ab}\otimes u_{ai}u_{bj}$$

This action only depends on the projective version $PG$, generated by the variables $u_{ai}u_{bj}$, and there is a biggest $G$ acting on $X$, see \cite{ba1}, \cite{bgo}, \cite{go1}. We call $PG$ the projective quantum isometry group of $X$. The following result is known, from \cite{ba1}, \cite{bgo}:

\begin{theorem}
\label{thm:projiso}
The projective quantum isometry groups of the projective spaces
$$P^{N-1}_\mathbb R\subset P^{N-1}_\mathbb C\subset P^{N-1}_+$$
are respectively $PO_N\subset PU_N\subset PO_N^+$, with $PU_N=PO_N^*$.
\end{theorem}

\begin{proof}
In the free case, the fact that $O_N^+$ acts is a direct computation, which will be done in a more general setting in section 4. In the real and complex cases, this is from \cite{ba1}, the idea being that the relations $p_{ab}=p_{ba}$ force $G\subset O_N$, and that the rank one relations $p_{ab}p_{cd}=p_{ad}p_{cb}$ force $G\subset O_N^*$.
\end{proof}

\section{Products and Segre}

Given two compact quantum spaces $X,Y$, their product is given by $C(X\times Y)=C(X)\otimes_{max}C(Y)$, the maximal tensor product, with the variables coming from $X$ commuting with those coming from $Y$. Classically, the Segre embedding is the map $P^{M-1}\times P^{N-1}\to P^{MN-1}$ given by $(x,y)\to x\otimes y$, which in terms of projections reads $(p,q)\to p\otimes q$. In the free setting, we have:

\begin{definition}
Given closed subspaces $X\subset P^{M-1}_+$ and $Y\subset P^{N-1}_+$, the free Segre map is
$$\sigma:X\times Y\to P^{MN-1}_+\quad,\quad
P_{(ia),(jb)}=p_{ij}q_{ab}$$
where $p,q$ are the projective coordinates of $X,Y$.
\end{definition}

Let us check that this is well-defined. The matrix $P=p\otimes q$ is self-adjoint, and since the entries of $p$ commute with those of $q$, we have $P^2=p^2\otimes q^2=P$, and $Tr(P)=Tr(p)Tr(q)=1$. Observe that in the free product setting, with the entries of $p,q$ not commuting, $P$ is no longer a projection, so the product $\times$ is the correct operation here.

\bigskip

We have the following result, which is the free analogue of the basic properties of the Segre embedding, and which works more generally for projections of arbitrary trace:

\begin{proposition}
\label{prop:segre}
Given projections $p\in M_M(A)$, $q\in M_N(A)$ of traces $L,K\geq1$, with commuting entries, $P=p\otimes q$ is a projection of trace $LK$, and we can recover $p,q$ from $P$ via
$$\sum_aP_{(ia),(ja)}=Kp_{ij}\quad,\quad
\sum_iP_{(ia),(ib)}=Lq_{ab}$$
so the Segre maps are embeddings of compact quantum spaces.
\end{proposition}

\begin{proof}
The projection property follows as above, and the trace is $Tr(p)Tr(q)=LK$. The partial trace formulae are clear, and show that the corresponding morphism of algebras is surjective, which gives the embedding property.
\end{proof}

Classically, the image of the Segre embedding is the space of rank one matrices, cut out by the $2\times2$ minors. In the free case, the minor relations are of commutation type, as the rank one relations from the proof of Theorem \ref{thm:projiso}, which reduce the quantum symmetries to the half-liberated level. However, we can use instead the partial traces, as follows:

\begin{proposition}
\label{prop:segreimage}
The image of the free Segre embedding $P^{M-1}_+\times P^{N-1}_+\subset P^{MN-1}_+$ is the closed subspace given by the relations
$$P_{(ia),(jb)}=p_{ij}q_{ab}\quad,\quad p_{ij}q_{ab}=q_{ab}p_{ij}\quad,\quad p^2=p\quad,\quad q^2=q$$
where $p_{ij}=\sum_aP_{(ia),(ja)}$ and $q_{ab}=\sum_iP_{(ia),(ib)}$ are the partial traces of $P$.
\end{proposition}

\begin{proof}
These relations hold on the image, by Proposition \ref{prop:segre}. Conversely, in the quotient $Q$ of $C(P^{MN-1}_+)$ by these relations, $p$ and $q$ are projections with commuting entries, and with $Tr(p)=Tr(q)=Tr(P)=1$, so we have a morphism $C(P^{M-1}_+)\otimes C(P^{N-1}_+)\to Q$ mapping the projective coordinates to the variables $p_{ij},q_{ab}$. By Proposition \ref{prop:segre} and by $P=p\otimes q$, this morphism is inverse to the one induced by the Segre map, and this gives the result.
\end{proof}

Here the quadratic relations $P=p\otimes q$ play the role of the $2\times2$ minors, the free setting forcing us to state them via the partial traces.

\bigskip

Regarding integration, the product integration on $X\times Y$ corresponds to independent copies, with commuting legs. For instance, for the free Grassmannians, with the integration from section 7, the Segre coordinate $P_{(11),(11)}=p_{11}q_{11}$ has as law the classical multiplicative convolution of the laws of $p_{11}$ and $q_{11}$. Thus, by Theorem \ref{thm:poisson} below, with $M,N\to\infty$ the variable $MNP_{(11),(11)}$ converges to the product of 2 independent, commuting, free Poisson variables, of parameters $L,K$.

\bigskip

Classically, the free Segre maps between Grassmannians correspond to $(E,F)\to E\otimes F$, for subspaces $E\subset\mathbb C^M$, $F\subset\mathbb C^N$ of dimensions $L,K$, and the image is the space of decomposable subspaces of $\mathbb C^M\otimes\mathbb C^N$, of dimension $LK$.

\bigskip

Regarding symmetries, if $G\subset O_M^+$ acts on $X$ and $H\subset O_N^+$ acts on $Y$, then $G\times H$ acts on $X\times Y$, and the Segre map is equivariant, with $G\times H$ acting on $P^{MN-1}_+$ via the tensor product corepresentation $u\otimes v$. Summarizing, we have:

\begin{theorem}
\label{thm:segre}
The free Segre maps $Gr^+_{LM}\times Gr^+_{KN}\to Gr^+_{LK,MN}$, and in particular
$$P^{M-1}_+\times P^{N-1}_+\subset P^{MN-1}_+$$
with $K,L\geq1$, are embeddings, which are equivariant with respect to the morphism $O_M^+\times O_N^+\to O_{MN}^+$ given by $w=u\otimes v$.
\end{theorem}

\begin{proof}
This follows from the above, with the Grassmannians $Gr^+_{LN}$ being given by the projections of trace $L$, see section 4. As for the equivariance, with $w=u\otimes v$ we have $w_{(ia),(kc)}w_{(jb),(ld)}=u_{ik}u_{jl}\otimes v_{ac}v_{bd}$, by commutativity of the two legs, and this gives the result.
\end{proof}

\section{Veronese and quadrics}

Classically, the Veronese embedding of degree $d$ is given by $x\to x^{\otimes d}$, which in terms of projections reads $p\to p^{\otimes d}$. In the free case $p^{\otimes d}$ is not a projection, due to the noncommutativity. However, the sphere picture suggests the correct definition. Indeed, for a point of the free real sphere, the variables $x_I=z_{i_1}\ldots z_{i_d}$ satisfy
$$\sum_Ix_Ix_I^*=\sum_Iz_{i_1}\ldots z_{i_d}z_{i_d}\ldots z_{i_1}=1$$
by cancelling the inner letters in pairs, and the projective version $P_{IJ}=x_Ix_J^*$ is given by the words $i_1\ldots i_dj_d\ldots j_1$, which can be written in terms of the projective coordinates $p_{ij}=z_iz_j$, by grouping the letters in pairs. We are led in this way to:

\begin{definition}
The free Veronese map of degree $d$ is the map $\nu_d:P^{N-1}_+\to P^{N^d-1}_+$ given by
$$P_{IJ}=p_{w_1w_2}p_{w_3w_4}\ldots p_{w_{2d-1}w_{2d}}$$
where $w=i_1\ldots i_dj_d\ldots j_1$ is the word associated to $I,J$.
\end{definition}

Observe that this word is palindromic precisely when $I=J$. For instance at $d=2$ we have $P_{(i_1i_2),(j_1j_2)}=p_{i_1i_2}p_{j_2j_1}$, and at $d=3$ we have $P_{(i_1i_2i_3),(j_1j_2j_3)}=p_{i_1i_2}p_{i_3j_3}p_{j_2j_1}$. We have the following result:

\begin{theorem}
\label{thm:veronese}
The free Veronese map $\nu_d:P^{N-1}_+\to P^{N^d-1}_+$ is well-defined, and is an embedding.
\end{theorem}

\begin{proof}
The adjoint of $P_{JI}$ is the product over the reversed word, which is again $P_{IJ}$, so $P=P^*$. Regarding now $P^2$, in the product $P_{IK}P_{KJ}$ the inner letters $k_1\ldots k_d$ appear in the middle, as $\ldots p_{k_2k_1}p_{k_1k_2}\ldots$, and summing over $k_1$, then $k_2$, and so on, by using $p^2=p$ and $Tr(p)=1$, these middle factors telescope:
$$\sum_{k_1}p_{k_2k_1}p_{k_1k_2}=p_{k_2k_2}\quad,\quad\sum_{k_2}p_{k_2k_2}=1\quad,\quad\ldots$$

At $d$ even we obtain in this way $P_{IJ}$, and at $d$ odd the last step is $\sum_{k_d}p_{i_dk_d}p_{k_dj_d}=p_{i_dj_d}$, and we obtain $P_{IJ}$ as well. The same telescoping gives $Tr(P)=1$. Finally, by using $Tr(p)=1$ we have $p_{ab}=\sum_{t_2\ldots t_d}p_{ab}p_{t_2t_2}\ldots p_{t_dt_d}$, and each term of this sum is a coordinate $P_{IJ}$, so $\nu_d$ is an embedding.
\end{proof}

As an illustration, the program {\tt verdi.py} checks the projection properties of the free Veronese maps, of the free Segre maps, and of the embeddings in Theorem \ref{thm:plucker} below, in random matrix models, where $p_{ij}=\sum_{r\leq L}c_{ri}c_{rj}A_iA_j$, with $c$ being a real $L\times N$ matrix with orthonormal rows, and with $A_1,\ldots,A_N$ being random self-adjoint unitaries.

\bigskip

Observe that classically, the Veronese embedding lands in the symmetric tensors, $P^{M-1}$ with $M=\binom{N+d-1}{d}$. In the free case there is no symmetrization, and we use the full tensor power. On the real points $p=zz^t$ we obtain $P_{IJ}=y_Iy_J$ with $y=z^{\otimes d}$, which is the usual Veronese map, so for instance $\nu_2(P^1_\mathbb R)$ is the conic $xz=y^2$, and $\nu_3(P^1_\mathbb R)$ is the twisted cubic, cut out by the $2\times2$ minors of a $2\times3$ matrix. On the complex points $p=zz^*$, however, we obtain $P_{IJ}=y_I\bar{y}_J$ with $y=z\otimes\bar{z}\otimes z\otimes\ldots$, with $d$ factors, so we have a real algebraic map, and for instance $\nu_2$ maps $P^1_+=P^1_\mathbb C=S^2$ onto the Bloch sphere of the rank one projections in $M_2(\mathbb C)$, which spans $P^3_\mathbb C$. The first genuinely free examples are the images $\nu_d(P^2_+)$, which are free analogues of the Veronese surfaces. In the free case such minor relations are of commutation type, and the images can be described instead via partial traces, as for the Segre embedding.

\bigskip

Let us discuss now hypersurfaces. Since in the sphere picture $p_{ij}=z_iz_j$, the linear relations between projective coordinates correspond to quadratic relations between affine coordinates. We are led to:

\begin{definition}
A free hypersurface of degree $d$ in $P^{N-1}_+$ is a closed subspace given by a relation
$$\sum_{a,b}c_{ab}\,p_{a_1b_1}\ldots p_{a_db_d}=0$$
with $c_{ab}\in\mathbb C$, and the free quadrics are the free hypersurfaces of degree $1$.
\end{definition}

As an example, the free conic $p_{11}+p_{22}=p_{33}$ in $P^2_+$ corresponds to the conic $x^2+y^2=z^2$ on the real points, while on the complex points it gives $|z_1|^2+|z_2|^2=|z_3|^2$, which is a real 3-sphere. The holomorphic conic $z_1^2+z_2^2=z_3^2$ corresponds instead to $\sum_{ab}s_as_bp_{ab}^2=0$, with $s=(1,1,-1)$, and more generally, a classical hypersurface $f(z)=0$ of degree $d$ in $P^{N-1}_\mathbb C$ can be written as $|f(z)|^2=0$, which is a relation of degree $d$ in the variables $p_{ab}=z_a\bar{z}_b$. For instance the Fermat cubic $\sum_az_a^3=0$ corresponds to $\sum_{ab}p_{ab}^3=0$. As in the classical case, the Veronese maps turn hypersurfaces into hyperplane sections, linear in the projective coordinates $P_{IJ}$:

\begin{theorem}
\label{thm:hyper}
The free hypersurfaces of degree $d$ in $P^{N-1}_+$ are exactly the preimages, under the free Veronese map $\nu_d$, of the hyperplane sections $\sum_{IJ}c_{IJ}P_{IJ}=0$ of $P^{N^d-1}_+$.
\end{theorem}

\begin{proof}
Any product $p_{a_1b_1}\ldots p_{a_db_d}$ is a product over the consecutive pairs of the word $a_1b_1\ldots a_db_d$, of length $2d$, and any word of length $2d$ is of the form $i_1\ldots i_dj_d\ldots j_1$. Thus the products of $d$ projective coordinates are exactly the coordinates $P_{IJ}$, and this gives the result.
\end{proof}

\section{Grassmannians}

In analogy with the free projective spaces, given by the projections of trace 1, we formulate:

\begin{definition}
Given integers $0\leq L\leq N$, the free Grassmannian $Gr^+_{LN}$ is given by
$$C(Gr^+_{LN})=C^*\left((p_{ij})_{i,j=1,\ldots,N}\Big|p=p^*=p^2,Tr(p)=L\right)$$
with the generators $p_{ij}$ being called projective coordinates.
\end{definition}

As for the free projective spaces, these algebras are well-defined, because the entries of a projection have norm $\leq1$, and the trace condition is not a consequence of the other relations, as we will see in Proposition \ref{prop:brown} below. Here are some basic properties of these spaces, which are all elementary:

\begin{proposition}
The free Grassmannians have the following properties:
\begin{enumerate}
\item $Gr^+_{0N}$ and $Gr^+_{NN}$ are points, and $Gr^+_{1N}=P^{N-1}_+$.

\item We have an isomorphism $Gr^+_{LN}\simeq Gr^+_{N-L,N}$, given by $p\to1-p$.

\item The classical version of $Gr^+_{LN}$ is the complex Grassmannian $Gr^\mathbb C_{LN}$.

\item The real Grassmannian $Gr^\mathbb R_{LN}$ appears by imposing $p=\bar{p}$, and commutativity.
\end{enumerate}
\end{proposition}

\begin{proof}
(1) At $L=0$ Lemma \ref{lem:trace} gives $p=0$, and at $L=N$ it gives $1-p=0$.

(2) This is clear, because $1-p$ is a projection, of trace $N-L$.

(3,4) These are clear, the characters corresponding to the projections of trace $L$.
\end{proof}

Let us mention that, in the absence of the trace condition, the algebra generated by the entries of a projection $p\in M_N(A)$ is the noncommutative Grassmannian of Brown \cite{bro}, studied by McClanahan \cite{mcc}. Thus $Gr^+_{LN}$ is the quotient of this algebra by $Tr(p)=L$. Regarding the relation between the 2 objects, we have the following result:

\begin{proposition}
\label{prop:brown}
The trace $Tr(p)$ is not central in the noncommutative Grassmannian of Brown, for any $N\geq2$. Also, the quotient of this algebra by $Tr(p)=\lambda$ is nonzero for $\lambda=r/N$, with $r\in\{0,1,\ldots,N^2\}$.
\end{proposition}

\begin{proof}
We use models with $A=M_K(\mathbb C)$, where a projection $p\in M_N(A)$ is a projection $P\in M_N(\mathbb C)\otimes M_K(\mathbb C)$, and $Tr(p)$ is the partial trace of $P$ over the first factor. For $P$ the projection onto a unit vector $v=\sum_ie_i\otimes v_i$ we have $p_{ij}=v_iv_j^*$. With $K=2$, $v_1=e_1/\sqrt{2}$, $v_2=(e_1+e_2)/2$ and $v_i=0$ for $i\geq3$, the element $Tr(p)=v_1v_1^*+v_2v_2^*$ does not commute with $p_{11}=v_1v_1^*$. For the second assertion, take $K=N$, and consider the Weyl matrices $W_{ab}=X^aZ^b$, with $a,b\in\mathbb Z_N$, where $X$ is the cyclic shift and $Z=diag(1,w,\ldots,w^{N-1})$, with $w=e^{2\pi i/N}$. The vectors $v_{ab}=N^{-1/2}\sum_ie_i\otimes W_{ab}e_i$ form an orthonormal basis of $\mathbb C^N\otimes\mathbb C^N$, because $Tr(W_{cd}^*W_{ab})=N\delta_{ac}\delta_{bd}$, and each of them has partial trace $W_{ab}W_{ab}^*/N=1/N$. Thus the projection onto the span of $r$ such vectors has partial trace $r/N$.
\end{proof}

Thus, besides the free Grassmannians, there are free spaces of projections with non-integer trace, as is the case for the scalar sums of projections from \cite{krs}, which appear here as diagonal examples, $p=diag(q_1,\ldots,q_N)$.

\bigskip

Brown also introduced in \cite{bro} the algebra generated by the entries of a unitary $u\in M_N(A)$, which is a dual group in the sense of Voiculescu, and whose probabilistic aspects were investigated by Franz and his collaborators, see \cite{bcf}. Also, Voiculescu extended in \cite{voi} his fully matricial function theory over a Banach algebra $B$ from the affine case to a Grassmannian completion, obtained by tying together certain Grassmannians for the matrix algebras $M_n(B)$, as a noncommutative analogue of the Riemann sphere $P^1_\mathbb C$, which is our $P^1_+$ as well, and this framework was extended to Grassmannians and flag manifolds by Ito \cite{ito}.

\bigskip

Classically, the Grassmannian is embedded into a projective space via the Pl\"ucker embedding, which uses the determinant. In the free case there is no determinant, but we can use instead the following construction, inspired by the Veronese map:

\begin{theorem}
\label{thm:plucker}
We have embeddings $Gr^+_{LN}\subset P^{N^2-1}_+$ and $Gr^+_{LN}\subset S^{N^2-1}_{\mathbb C,+}$, given by
$$P_{(i_1i_2),(j_1j_2)}=\frac{1}{L}\,p_{i_1i_2}p_{j_2j_1}\quad,\quad
x_{ij}=\frac{p_{ij}-\frac{L}{N}\delta_{ij}}{\sqrt{L(N-L)/N}}$$
the first one being valid at $L\geq1$, and the second one at $1\leq L\leq N-1$, with the coordinates $x_{ij}$ spanning the first nonconstant eigenspace of the Brownian Laplacian, in the orbit model.
\end{theorem}

\begin{proof}
For the first embedding, $P$ is self-adjoint, and the computation in the proof of Theorem \ref{thm:veronese} gives $P^2=P$, the middle factor being now $Tr(p^2)/L=1$, and $Tr(P)=Tr(p^2)/L=1$. We can recover $p$ via $p_{ij}=\sum_kP_{(ij),(kk)}$. For the second embedding, with $c=L/N$ we have:
$$\sum_{ij}(p_{ij}-c\delta_{ij})(p_{ij}-c\delta_{ij})^*
=Tr((p-c)^2)
=L-2cL+c^2N
=\frac{L(N-L)}{N}$$

The same computation works for $\sum_{ij}x_{ij}^*x_{ij}$, and this gives the result. Regarding the eigenspace claim, we will see in sections 7-8 that the variables $x_{ij}$ span the first nonconstant eigenspace, as for the minimal embeddings in the classical case, see Takahashi \cite{tak}.
\end{proof}

Classically, the Grassmannian $Gr_{2,4}$ is the Klein quadric in $P^5$, via the Pl\"ucker embedding, which uses the exterior square of $\mathbb C^4$. In the free case there is no exterior square, and our embedding $Gr^+_{2,4}\subset P^{15}_+$ uses the full tensor square, with the image being given, as in Proposition \ref{prop:segreimage}, by the relations $P_{(i_1i_2),(j_1j_2)}=\frac{1}{L}p_{i_1i_2}p_{j_2j_1}$, where $p_{ij}=\sum_kP_{(ij),(kk)}$ is a projection of trace $L$. Finding a free analogue of the Klein quadric, and of the Pl\"ucker relations, would be very interesting, the toric deformations of the Pl\"ucker relations being discussed in \cite{cls}.

\bigskip

Regarding symmetries, a closed subgroup $G\subset O_N^+$ acts on a closed subspace $X\subset Gr^+_{LN}$ when the formula $\Phi(p_{ij})=\sum_{ab}p_{ab}\otimes u_{ai}u_{bj}$ defines a morphism, and we define the projective quantum isometry group as in section 1. In order to deal with the real case, we will need:

\begin{lemma}
\label{lem:indep}
For $1\leq L\leq N-1$, the coordinates $p_{ab}$ with $a\leq b$ are linearly independent, as functions on the real Grassmannian $Gr^\mathbb R_{LN}$.
\end{lemma}

\begin{proof}
Assume $\sum_{a\leq b}c_{ab}p_{ab}=0$. With $C$ symmetric, this reads $Tr(Cp)=0$ for any real projection $p$ of rank $L$. Given unit vectors $x,y\in\mathbb R^N$, we can find a subspace $F\perp x,y$ of dimension $L-1$, because $N-2\geq L-1$. With $p,p'$ being the projections onto $F+\mathbb Rx$ and $F+\mathbb Ry$, we obtain $Tr(C(xx^t-yy^t))=0$, so $x^tCx$ is constant on the sphere, and so $C=c1$. Now $Tr(Cp)=cL=0$ gives $c=0$.
\end{proof}

In the real and free cases, we have the following result:

\begin{proposition}
\label{prop:grassiso}
The projective quantum isometry groups of the real and free Grassmannians $Gr^\mathbb R_{LN}\subset Gr^+_{LN}$ with $1\leq L\leq N-1$ are respectively $PO_N\subset PO_N^+$.
\end{proposition}

\begin{proof}
In the free case, $P=(1\otimes u^t)(p\otimes1)(1\otimes u)$ is a projection of trace $L$, by using $uu^t=u^tu=1$, so $O_N^+$ acts. In the real case, assume that $G\subset O_N^+$ acts. Since $p_{ij}=p_{ji}$, we have $\Phi(p_{ij})=\Phi(p_{ji})$, and with $w_{ij,ab}=u_{ai}u_{bj}$ we obtain:
$$\sum_{a<b}p_{ab}\otimes(w_{ij,ab}+w_{ij,ba}-w_{ji,ab}-w_{ji,ba})+\sum_ap_{aa}\otimes(w_{ij,aa}-w_{ji,aa})=0$$

By Lemma \ref{lem:indep} all the coefficients vanish, and by doubling the diagonal ones, we obtain $w_{ij,ab}+w_{ij,ba}=w_{ji,ab}+w_{ji,ba}$, for any $a,b$. Now by applying the antipode, which maps $w_{ij,ab}\to w_{ba,ji}$, and relabeling, we obtain:
$$w_{ji,ba}+w_{ij,ba}=w_{ji,ab}+w_{ij,ab}$$

By comparing with the original relation, we obtain $w_{ij,ab}=w_{ji,ba}$, that is, $u_{ai}u_{bj}=u_{bj}u_{ai}$. Thus $G\subset O_N$, and so $PG\subset PO_N$, and the converse inclusion is clear.
\end{proof}

In the complex case, the half-liberation $O_N^*$ acts on $Gr^\mathbb C_{LN}$, because via the standard matrix model $u_{ij}\to(^{\ 0\ }_{\bar{U}_{ij}}{\ }^{U_{ij}}_{\ 0\ })$, with $U\in U_N$, the coaction becomes $p\to(U^tp\bar{U},U^*pU)$. At $L=1$ the rank one relations force $PG\subset PU_N$, as explained in section 1, but at $L\geq2$ these relations are not available any longer. We will use instead the following relations:

\begin{lemma}
\label{lem:isotropic}
Assume that $G\subset O_N^+$ acts on $Gr^\mathbb C_{LN}$, with $1\leq L\leq N-1$. Then for any $x,y\in\mathbb C^N$ with $\sum_ax_ay_a=0$, the variables $z_i=\sum_ax_au_{ai}$ and $t_i=\sum_ay_au_{ai}$ satisfy $[z_it_j,z_kt_l]=0$, for any $i,j,k,l$.
\end{lemma}

\begin{proof}
Since $C(Gr^\mathbb C_{LN})$ is commutative, we have $[\Phi(p_{ij}),\Phi(p_{kl})]=0$, and by evaluating at a projection $q$ of rank $L$, we obtain that the elements $Y_{ij}=\sum_{ab}q_{ab}u_{ai}u_{bj}$ commute. Now fix orthonormal vectors $\xi_2,\ldots,\xi_L$, let $K$ be their orthogonal complement, and take $q=\xi\xi^*+\sum_{r\geq2}\xi_r\xi_r^*$, with $\xi\in K$ of norm 1. With $A_{ij}=\sum_{ab}\xi_a\bar{\xi}_bu_{ai}u_{bj}$ and $C_{ij}$ being the remaining part of $Y_{ij}$, by homogeneity we obtain, for any $\xi\in K$:
$$[A_{ij},A_{kl}]+|\xi|^2([A_{ij},C_{kl}]+[C_{ij},A_{kl}])+|\xi|^4[C_{ij},C_{kl}]=0$$

Let $f_1,f_2,\ldots$ be an orthonormal basis of $K$, and write $\xi=\sum_s\lambda_sf_s$. The coefficient of $\lambda_1^2\bar{\lambda}_2^2$ comes only from the first term, and gives the relation in the statement, with $x=f_1$ and $y=\bar{f}_2$. Since any $x,y$ with $\sum_ax_ay_a=0$ appear in this way, up to scaling, because $N\geq L+1$, this gives the result.
\end{proof}

In tensor notation, with $c=x\otimes y\otimes x\otimes y$, and with $T$ being the flip $e_a\otimes e_b\otimes e_c\otimes e_d\to e_c\otimes e_d\otimes e_a\otimes e_b$, these relations read $c^tu^{\otimes4}=c^tu^{\otimes4}T$. Consider the space $V\subset(\mathbb C^N)^{\otimes4}$ of vectors which are symmetric in the legs $1,3$ and in the legs $2,4$, and whose contractions between an odd leg and an even leg vanish, and the space $W$ defined similarly, with antisymmetric instead of symmetric. Via the above matrix model, these are the trace-free parts of $S^2\mathbb C^N\otimes S^2\overline{\mathbb C}^N$ and $\Lambda^2\mathbb C^N\otimes\Lambda^2\overline{\mathbb C}^N$, which are irreducible representations of $U_N$, with highest weights $(2,0,\ldots,0,-2)$ and $(1,1,0,\ldots,0,-1,-1)$, non-isomorphic for $N\geq4$, see Fulton and Harris \cite{fha}. The vector $e_1\otimes e_N\otimes e_1\otimes e_N$ is one of the vectors $c$, and is a highest weight vector, so the vectors $c$ span $V$. We have:

\begin{proposition}
\label{prop:tannakian}
Assume that $O_N^*\subset G\subset O_N^+$ acts on $Gr^\mathbb C_{LN}$, with $1\leq L\leq N-1$ and $N\geq4$. Then $End(u^{\otimes4})$ contains the orthogonal projection onto $V$, or the one onto $V\oplus W$.
\end{proposition}

\begin{proof}
Let $v=u^{\otimes4}$. By Lemma \ref{lem:isotropic}, the matrix coefficients $c^tva$ vanish, for $c\in V$ and $Ta=-a$. By Peter-Weyl, such a coefficient vanishes precisely when $c^tXa=0$ for any $X$ in the commutant $End(v)'$. Since $\bar{V}=V$, this says that $Xa\perp V$, and by applying this to $X^*$, we obtain $XV\perp a$. Thus the space $\widetilde{V}=End(v)'V$, whose projection belongs to $End(v)''=End(v)$, is orthogonal to $\ker(T+1)$, and so is contained in $\ker(T-1)$. Also, since $End(v)$ contains the Jones-Wenzl projection $P_4$, we have $\widetilde{V}\subset r_4$. Now a vector fixed by $T$ decomposes into its symmetric and antisymmetric parts with respect to the legs $1,3$, which are the same as those with respect to the legs $2,4$, and the conditions defining $r_4$, together with the invariance under $T$, kill all the odd-even contractions, so $r_4\cap\ker(T-1)=V\oplus W$. Finally, since $O_N^*\subset G$, the space $\widetilde{V}$ is invariant under $U_N$, and since $V,W$ are irreducible and non-isomorphic, we obtain $\widetilde{V}=V$ or $\widetilde{V}=V\oplus W$.
\end{proof}

These projections can be written in terms of the 24 walled Brauer diagrams on 8 points, which span $End_{O_N^*}(u^{\otimes4})$. With $D_k$ being the set of such diagrams having $k$ caps on top, and $k$ cups at the bottom, the program {\tt mascagni.py} finds the following formula, and a similar one for the projection onto $W$, with denominators $N-2$ and $(N-1)(N-2)$:
$$P_V=\frac{1}{4}\sum_{\pi\in D_0}T_\pi-\frac{1}{4(N+2)}\sum_{\pi\in D_1}T_\pi+\frac{1}{2(N+1)(N+2)}\sum_{\pi\in D_2}T_\pi$$

The same program then finds, starting from the Temperley-Lieb maps and from $P_V$, respectively from $P_V+P_W$, explicit expressions of $\sigma\otimes1$, with $\sigma:e_a\otimes e_b\otimes e_c\to e_c\otimes e_b\otimes e_a$ being the half-liberated crossing, as linear combinations, with coefficients in $\mathbb Q(N)$, of products, adjoints and rotations of these maps. The denominators of these coefficients vanish only at $N\in\{0,\pm1,\pm2,-3\}$ and at $N^2\in\{6,7,8\}$, as checked by exact factorization, the program also checking that the residual vanishes identically. We obtain in this way the following result, which completes Theorem \ref{thm:projiso}:

\begin{theorem}
\label{thm:grassiso}
The projective quantum isometry groups of the real, complex and free Grassmannians
$$Gr^\mathbb R_{LN}\subset Gr^\mathbb C_{LN}\subset Gr^+_{LN}$$
with $1\leq L\leq N-1$ are respectively $PO_N\subset PU_N\subset PO_N^+$.
\end{theorem}

\begin{proof}
The real and free cases are from Proposition \ref{prop:grassiso}. In the complex case, $PU_N=PO_N^*$ acts, as explained above. Conversely, at $L\in\{1,N-1\}$ the result follows from Theorem \ref{thm:projiso}, by using $p\to1-p$, so assume $2\leq L\leq N-2$, and so $N\geq4$. Let $G\supset O_N^*$ be the biggest quantum group acting. By Proposition \ref{prop:tannakian}, $End(u^{\otimes4})$ contains $P_V$ or $P_V+P_W$. Since the tensor category of $G$ contains the Temperley-Lieb maps, and is stable under products, adjoints and rotations, the above computation gives $\sigma\otimes1\in End(u^{\otimes4})$. By closing the last string we obtain $\sigma\in End(u^{\otimes3})$, so $G\subset O_N^*$, and so $PG\subset PO_N^*=PU_N$.
\end{proof}

\section{Flag manifolds}

Classically, the flag manifolds are the spaces of increasing sequences of subspaces $E_1\subset\ldots\subset E_k=\mathbb C^N$ of given dimensions. In terms of the projections onto the successive orthogonal complements, we are led to:

\begin{definition}
Given a composition $d=(d_1,\ldots,d_k)$ of $N$, the free flag manifold $Fl^+_d$ is given by
$$C(Fl^+_d)=C^*\left((q^s_{ij})_{s\leq k,\,i,j\leq N}\Big|q^s=(q^s)^*=(q^s)^2,\ \sum_sq^s=1,\ Tr(q^s)=d_s\right)$$
and the full free flag manifold is $Fl^+_N=Fl^+_{(1,\ldots,1)}$.
\end{definition}

Observe that we have $q^sq^t=0$ for $s\neq t$, as for any family of projections summing up to 1, so the partial sums $q^1+\ldots+q^s$ are projections, of traces $d_1+\ldots+d_s$. Classically, the complex flag manifold is $Fl^\mathbb C_d=U_N/(U_{d_1}\times\ldots\times U_{d_k})$, of real dimension $N^2-\sum_sd_s^2$, with the full flag manifold $Fl^\mathbb C_N=U_N/\mathbb T^N$ having real dimension $N(N-1)$. Also, given distinct real numbers $\lambda_1,\ldots,\lambda_k$, the flags correspond to the self-adjoint matrices $x=\sum_s\lambda_sq^s$, which form the isospectral manifold, namely the adjoint orbit $\{UDU^*|U\in U_N\}$ of the matrix $D=diag(\lambda_1^{(d_1)},\ldots,\lambda_k^{(d_k)})$, with $\lambda^{(d)}$ meaning $\lambda$ repeated $d$ times, consisting of the Hermitian matrices with this prescribed spectrum. In the free setting, we have:

\begin{proposition}
\label{prop:isospectral}
The free flag manifold $Fl^+_d$ is the free isospectral manifold
$$C(Fl^+_d)=C^*\left((x_{ij})_{i,j=1,\ldots,N}\Big|x=x^*,\ \prod_s(x-\lambda_s)=0,\ Tr(x^m)=\sum_sd_s\lambda_s^m,\forall m<k\right)$$
and we have maps $Fl^+_d\to Gr^+_{d_1+\ldots+d_s,N}$, given by $p=q^1+\ldots+q^s$, and $Fl^+_d\to Fl^+_{d'}$, for any coarsening $d'$ of $d$. Also, we have $Fl^+_{(L,N-L)}=Gr^+_{LN}$.
\end{proposition}

\begin{proof}
Given $x$ as above, the spectral projections $q^s=\prod_{t\neq s}\frac{x-\lambda_t}{\lambda_s-\lambda_t}$ are polynomials in $x$, they are projections summing up to 1, and the trace conditions give, via the Vandermonde matrix, $Tr(q^s)=d_s$. Conversely, $x=\sum_s\lambda_sq^s$. The other assertions are clear.
\end{proof}

Observe that the trace conditions are needed here, because $\prod_s(x-\lambda_s)=0$ alone only says that the spectrum of $x$ is contained in $\{\lambda_1,\ldots,\lambda_k\}$, and allows for instance $x=\lambda_11$, the multiplicities $d_s$ being fixed by the trace conditions. Observe also that, classically, the isospectral manifolds are the coadjoint orbits of $U_N$, with their Kirillov-Kostant symplectic structure. For instance at $k=2$ we have $x=\lambda_1q^1+\lambda_2(1-q^1)$, and we are back to the Grassmannians, with $p=q^1$ being a projection of trace $d_1$. The main source of free flags is the free orthogonal group $O_N^+$, as follows:

\begin{theorem}
\label{thm:flagmaps}
Let $B_1,\ldots,B_k$ be the consecutive blocks of $\{1,\ldots,N\}$ of sizes $d_1,\ldots,d_k$.
\begin{enumerate}
\item We have a morphism $C(Fl^+_d)\to C(O_N^+)$ given by $q^s_{ij}\to\sum_{r\in B_s}u_{ri}u_{rj}$.

\item The image is invariant under the left action of the block-diagonal subgroup of $O_N^+$, which is the free product $O_{d_1}^+\,\widehat{*}\,\ldots\,\widehat{*}\,O_{d_k}^+$, in the sense of Wang \cite{wa1}.

\item We have a coaction $C(Fl^+_d)\to C(Fl^+_d)\otimes C(O_N^+)$, given by $q^s_{ij}\to\sum_{ab}q^s_{ab}\otimes u_{ai}u_{bj}$.
\end{enumerate}
In particular, for the full flag, the orbit map $O_N^+\to Fl^+_N$ factors through the quotient $O_N^+/T_N^+$, taken with respect to the left action of the free torus $T_N^+$.
\end{theorem}

\begin{proof}
(1) With $Q^s_{ij}=\sum_{r\in B_s}u_{ri}u_{rj}$ the matrices $Q^s$ are self-adjoint. Using $uu^t=1$ we obtain $\sum_lQ^s_{il}Q^t_{lj}=\delta_{st}Q^s_{ij}$, using $u^tu=1$ we obtain $\sum_sQ^s=1$, and using $uu^t=1$ again we obtain $Tr(Q^s)=\sum_{r\in B_s}\sum_iu_{ri}^2=d_s$.

(2) The left action of the block-diagonal subgroup is $u_{ri}\to\sum_xg_{rx}\otimes u_{xi}$, with $x$ ranging over the block of $r$, and with $g$ being block-diagonal. Thus we have:
$$Q^s_{ij}\to\sum_{x,y\in B_s}\left(\sum_{r\in B_s}g_{rx}g_{ry}\right)\otimes u_{xi}u_{yj}=Q^s_{ij}$$

(3) With $P^s=(1\otimes u^t)(q^s\otimes1)(1\otimes u)$, the relations follow from $uu^t=u^tu=1$. Finally, at $k=N$ the block-diagonal subgroup is $\widehat{\mathbb Z_2^{*N}}=T_N^+$.
\end{proof}

In the complex case, the same formulae with $u_{ri}u_{rj}$ replaced by $u_{ri}^*u_{rj}$ give morphisms $C(Fl^+_d)\to C(U_N^+)$. As for the free projective spaces, nothing new comes in this way:

\begin{theorem}
\label{thm:realcomplex}
The real and complex maps $\alpha_\mathbb R:C(Fl^+_d)\to C(O_N^+)$ and $\alpha_\mathbb C:C(Fl^+_d)\to C(U_N^+)$ satisfy
$$\int_{O_N^+}\alpha_\mathbb R(f)=\int_{U_N^+}\alpha_\mathbb C(f)$$
for any $f\in C(Fl^+_d)$. Moreover, the images of $\alpha_\mathbb R,\alpha_\mathbb C$ are isomorphic, so the real and complex orbit models of the free flag manifolds, and in particular of the free Grassmannians, coincide.
\end{theorem}

\begin{proof}
It is enough to check this on products of generators. In the real case, these become sums of terms $\int u_{i_1j_1}\ldots u_{i_{2k}j_{2k}}$, and in the complex case, sums of the same terms, with the colors $u^*,u,u^*,u,\ldots$ on the entries. By the Weingarten formula, these integrals involve the sets $NC_2(2k)$ and $\mathcal{NC}_2(\bullet\circ\bullet\circ\ldots)$, which are equal, because a noncrossing pairing always connects points of different parity. Since the Gram matrices, $G(\pi,\sigma)=N^{|\pi\vee\sigma|}$, are equal too, this gives the first assertion. For the second one, with $v$ being the fundamental corepresentation of $U_N^+$, the matrix $g_{ij}=(^{\ 0\ }_{v_{ij}^*}{\ }^{v_{ij}}_{\ 0\ })$ is orthogonal, by biunitarity, and the lower corner of $\sum_{r\in B_s}g_{ri}g_{rj}$ is $\sum_{r\in B_s}v_{ri}^*v_{rj}$. Taking this corner is multiplicative on the block-diagonal image, and we obtain a morphism $Im(\alpha_\mathbb R)\to Im(\alpha_\mathbb C)$, whose inverse comes from $v\to u$.
\end{proof}

Regarding symmetries, Theorem \ref{thm:grassiso} extends to the flag manifolds: if $G\subset O_N^+$ acts on $Fl^\mathbb C_d$, the variables $q^1_{ij}$ generate a copy of $C(Gr^\mathbb C_{d_1,N})$, stable under the coaction, so $PG\subset PU_N$, and $PU_N$ acts, as before. Thus, the projective quantum isometry group of $Fl^\mathbb C_d$ is $PU_N$, for any $d$ having at least 2 terms. Similarly, by using Proposition \ref{prop:grassiso}, the projective quantum isometry groups of the real and free flag manifolds are $PO_N$ and $PO_N^+$, the free case coming from Theorem \ref{thm:flagmaps} (3).

\bigskip

Let us discuss now the relation between the universal flag algebras and their orbit models. Let $\Phi$ be the coaction from Theorem \ref{thm:flagmaps} (3), and $Pol(Fl^+_d)\subset C(Fl^+_d)$ be the $*$-algebra generated by the coordinates $q^s_{ij}$. We have:

\begin{proposition}
\label{prop:ergodic}
The averaging map $E=(id\otimes h)\Phi$ satisfies $E(f)=h(\alpha_\mathbb R(f))1$, for any $f\in Pol(Fl^+_d)$. In particular, the coaction of $O_N^+$ on $Fl^+_d$ is ergodic, and $\alpha_\mathbb R$ is injective on $Pol(Fl^+_d)$.
\end{proposition}

\begin{proof}
By the Weingarten formula, $E$ maps a product $q^{s_1}_{i_1j_1}\ldots q^{s_k}_{i_kj_k}$ to a combination of contractions of the indices of $q^{s_1}_{a_1b_1}\ldots q^{s_k}_{a_kb_k}$ along noncrossing pairings. Each such pairing has an adjacent pair, contracting either $a_l,b_l$, which gives $Tr(q^s)=d_s$, or $b_l,a_{l+1}$, which gives $q^sq^t=\delta_{st}q^s$, so by induction all these contractions are scalars, and so is $E(f)$. By applying $\alpha_\mathbb R$, and using $\Delta\alpha_\mathbb R=(\alpha_\mathbb R\otimes id)\Phi$, this scalar is $h(\alpha_\mathbb R(f))$. Now if $\Phi(f)=f\otimes1$ then $f=E(f)$ is a scalar, which gives the ergodicity. Finally, assume $\alpha_\mathbb R(f)=0$, and write $\Phi(f)=\sum_ia_i\otimes w_i$, with $w_i\in Pol(O_N^+)$ linearly independent. Then $\sum_{ij}a_i^*a_jh(w_i^*w_j)=E(f^*f)=0$, and since the Gram matrix $(h(w_i^*w_j))$ is positive definite, we obtain $a_i=0$ for any $i$, and so $f=(id\otimes\varepsilon)\Phi(f)=0$.
\end{proof}

Thus, the orbit models are faithful at the polynomial level, the faithfulness of $C(Fl^+_d)\to C(O_N^+)$ at the full $C^*$-algebra level being an open question. Observe also that, due to the ergodicity, the coaction of $O_N^+$ on $Fl^+_d$ is transitive in the quantum sense, even though the classical coset description fails, as explained in section 7.

\bigskip

Let us discuss now the torus fixed points. The free torus $T_N^+\subset O_N^+$, given by $u=diag(g_1,\ldots,g_N)$ with $g_i$ being the standard generators of $\mathbb Z_2^{*N}$, acts on $Fl^+_d$ via $q^s_{ij}\to q^s_{ij}\otimes g_ig_j$. Given a coaction $\Phi:C(X)\to C(X)\otimes C(G)$, its fixed point space is the largest quotient of $C(X)$ on which $\Phi(f)=f\otimes1$. We have:

\begin{proposition}
The fixed points of $T_N^+\curvearrowright Fl^+_d$ are obtained by imposing the relations $q^s_{ij}=0$ for $i\neq j$, and the same holds for the classical torus $T_N=\mathbb Z_2^N\subset T_N^+$.
\end{proposition}

\begin{proof}
The fixed point relations are $q^s_{ij}\otimes(g_ig_j-1)=0$. At $i=j$ these are automatic, and at $i\neq j$, by applying a character of $\mathbb Z_2^N$ mapping $g_i\to1$, $g_j\to-1$, we obtain $q^s_{ij}=0$. Conversely, these latter relations make the coaction trivial.
\end{proof}

Thus, the fixed points are the diagonal flags, whose diagonal entries form a matrix of projections $m_{si}=q^s_{ii}$, with $\sum_sm_{si}=1$ and $\sum_im_{si}=d_s$. In the classical case these projections are scalars, and the diagonal flags are the coordinate flags, corresponding to the $N!/d_1!\ldots d_k!$ ways of distributing the $N$ coordinate axes among the blocks. For an arbitrary free flag, the diagonal entries still form a bistochastic type matrix, with entries $0\leq m_{si}\leq1$. In the image of $O_N^+$ we have $m_{si}=\sum_{r\in B_s}u_{ri}^2$, which for the full flag is $m_{si}=u_{si}^2$, the free analogue of the classical orthostochastic matrices, and we will come back to this in section 10, in relation with the mixing matrices from particle physics.

\bigskip

Regarding the Grassmannian case, the universal algebra of projections $q_1,\ldots,q_N$ summing up to a scalar $\lambda$ was studied by Kruglyak, Rabanovich and Samoilenko \cite{krs}, who determined the values of $\lambda$ for which such projections exist, and by Shulman \cite{shu}, who studied the analytic properties of these algebras. At $L=1$ the projections must be orthogonal, and we obtain $\mathbb C^N$, corresponding to the $N$ coordinate points of $P^{N-1}_+$, and the same happens at $L=N-1$, by complementation. At $2\leq L\leq N-2$ these algebras are noncommutative, because we have a quotient map onto the universal algebra generated by 2 projections $p,r$, which is the group algebra of the infinite dihedral group $\mathbb Z_2*\mathbb Z_2$, given by $q\to(p,1-p,r,1-r,1,\ldots,1,0,\ldots,0)$, with $L-2$ values 1 and $N-L-2$ values 0. Thus the torus fixed points of $Gr^+_{2,4}$ form an infinite noncommutative space, in contrast with the 6 coordinate points of the classical Grassmannian $Gr_{2,4}$. Classically, the coordinate flags are the centers of the Schubert cells, which are indexed by the cosets of $S_{d_1}\times\ldots\times S_{d_k}$ in the Weyl group $S_N$, their number being $N!/d_1!\ldots d_k!$, and a free analogue of the Schubert cells, and of the Schubert calculus, centered at the torus fixed points, with $S_N^+$ playing the role of the Weyl group, would be of great interest.

\bigskip

Getting back now to the torus fixed points, the above discussion leads to the following result, which includes Theorem B from the introduction:

\begin{theorem}
\label{thm:fixed}
The torus fixed points of the free flag manifolds are as follows:
\begin{enumerate}
\item For $Fl^+_N$ we obtain the quantum permutation group $S_N^+$.

\item For $Gr^+_{LN}$ we obtain the universal algebra of $N$ projections summing up to $L$.

\item For $Fl^+_d$ we obtain the universal algebra of the $k\times N$ matrices of projections $m$ with columns summing up to $1$, and with rows summing up to $d_1,\ldots,d_k$.
\end{enumerate}
Moreover, we have morphisms from these algebras to $C(S_N^+)$, given by $m_{si}\to\sum_{r\in B_s}u_{ri}$.
\end{theorem}

\begin{proof}
This follows from the above discussion. In the full flag case the matrix $m$ is $N\times N$, with entries being projections, and with rows and columns summing up to 1, so this is a magic unitary, and we obtain the algebra $C(S_N^+)$ of Wang \cite{wa2}. As for the last assertion, this comes from the fact that the entries in a column of a magic unitary are orthogonal projections, summing up to 1, and that the entries in a row sum up to 1.
\end{proof}

\section{Determinantal varieties}

Classically, the determinantal varieties are the spaces of matrices of rank $\leq r$, given by the vanishing of the $(r+1)\times(r+1)$ minors. In the free case the minors are not available, and at $r=1$ the $2\times2$ minor relations $x_{ab}x_{cd}=x_{ad}x_{cb}$ are of commutation type, and are not the right ones, as explained in section 2. However, we can use instead the classical incidence model, which resolves the singularities of the determinantal varieties, see \cite{har}:
$$\widetilde{D}_r=\left\{(X,\Lambda)\Big|Im(X)\subset\Lambda,\ \dim\Lambda=r\right\}$$

In terms of the projection $p$ onto $\Lambda$, the condition $Im(X)\subset\Lambda$ reads $pX=X$. We are led in this way to the following definition, where we normalize the matrices, as to obtain compact spaces:

\begin{definition}
The free determinantal variety $D^+_{rMN}$ is given by
$$C(D^+_{rMN})=C^*\left(x_{ij},p_{ab}\Big|x=\bar{x},\ \sum_{ij}x_{ij}^2=1,\ p=p^*=p^2,\ Tr(p)=r,\ px=x\right)$$
with $1\leq r\leq M$, with $x\in M_{M\times N}(A)$ having self-adjoint entries, and with $p\in M_M(A)$.
\end{definition}

Here the conditions $x=\bar{x}$ and $\sum_{ij}x_{ij}^2=1$ tell us that $x$ is a point of the free sphere $S^{MN-1}_{\mathbb R,+}$, and the conditions on $p$ tell us that $p\in Gr^+_{rM}$. Thus, we have maps as follows, with the second one being the free analogue of the projection onto the base of a vector bundle:
$$D^+_{rMN}\to S^{MN-1}_{\mathbb R,+}\quad,\quad
D^+_{rMN}\to Gr^+_{rM}$$

Classically, the commutative version of $D^+_{rMN}$ is a mixed incidence space, consisting of the real matrices $x$ of rank $\leq r$ with $\sum_{ij}x_{ij}^2=1$, together with complex subspaces $\Lambda\subset\mathbb C^M$ of dimension $r$ containing $Im(x)\otimes_\mathbb R\mathbb C$. The fiber over an $x$ of rank $s$ is the complex Grassmannian $Gr^\mathbb C_{r-s,M-s}$, which is a point at $s=r$. For $1\leq r\leq\min(M,N)$ the closure of the rank $r$ part consists of the pairs with $\Lambda=\bar{\Lambda}$, because at rank $r$ the subspace $\Lambda$ is the complexification of $Im(x)$. This closure is everything at $r=1$, and at $r=M\leq N$, but is proper for $2\leq r<M$, as shown by a rank one matrix with image $\mathbb Re_1$, paired with the span of $e_1,\ldots,e_{r-1},e_r+ie_{r+1}$. For $N<r\leq M$ there are no matrices of rank $r$. Observe also that, in the fully complex setting and before normalization, the incidence model has dimension $r(M-r)+rN$, with $r(M-r)$ coming from the choice of $\Lambda$, and $rN$ from the choice of $X$, and for $r\leq\min(M,N)$ this is $r(M+N-r)$, which is the dimension of the variety of matrices of rank $\leq r$, as it should. In the free case, we can regard the image of the map $D^+_{rMN}\to S^{MN-1}_{\mathbb R,+}$ as being the free spherical determinantal variety. Classically, at $r=1$, the spherical determinantal variety consists of the matrices $x_{ia}=z_iw_a$, with $z,w$ being unit vectors, and so is the spherical version of the image of the Segre embedding. Observe also that we could use as well the orthogonal complement of the kernel, leading to the relation $xq=x$, with $q\in Gr^+_{rN}$, for $r\leq N$.

\bigskip

Regarding symmetries, the quantum group $O_M^+\times O_N^+$ acts on $D^+_{rMN}$, as follows:
$$x_{ij}\to\sum_{kl}x_{kl}\otimes a_{ki}b_{lj}\quad,\quad
p_{ij}\to\sum_{kl}p_{kl}\otimes a_{ki}a_{lj}$$

Indeed, the relation $px=x$ is preserved, because $\sum_ma_{lm}a_{sm}=\delta_{ls}$, and the other relations are preserved too, as before. Observe that $O_N^+$ acts trivially on $p$, in agreement with the fact that classically, the right action of $O_N$ leaves invariant the image of $x$, and so the condition $Im(x)\subset\Lambda$. We can now formulate our main result regarding the free determinantal varieties, which relates them to the Segre embedding, and to the free partial isometries from \cite{ba1}:

\begin{theorem}
\label{thm:determinantal}
The free determinantal varieties have the following properties:
\begin{enumerate}
\item At $r=M$ we have $p=1$, and $D^+_{MMN}=S^{MN-1}_{\mathbb R,+}$.

\item At $r=1$ we have a map $S^{M-1}_{\mathbb R,+}\times S^{N-1}_{\mathbb R,+}\to D^+_{1MN}$, given by $x_{ia}=z_iw_a$ and $p_{ij}=z_iz_j$, whose projective version is the free Segre embedding.

\item For $r\leq N$, the rescaled free partial isometries $x=u/\sqrt{r}$, with $u\in O^{r+}_{MN}$, embed into $D^+_{rMN}$, with $p=uu^t$, as the subspace given by $xx^t=p/r$.
\end{enumerate}
\end{theorem}

\begin{proof}
(1) This follows from Lemma \ref{lem:trace}, applied to $1-p$.

(2) With $z,w$ having commuting entries, we have $\sum_{ia}x_{ia}^2=\sum_iz_i^2\sum_aw_a^2=1$, and $(px)_{ia}=\sum_kz_iz_kz_kw_a=z_iw_a$. The projective version is $x_{ia}x_{jb}=z_iz_j\otimes w_aw_b$, which is the free Segre map.

(3) Given $u\in O^{r+}_{MN}$, meaning $u=\bar{u}$ with $p=uu^t$ being a projection of trace $r$, we have $X=u-pu$ satisfying $XX^*=0$, so $pu=u$. Also, $\sum_{ij}u_{ij}^2=Tr(p)=r$, which gives the result.
\end{proof}

Let us discuss now the Hermitian case. Classically, by the spectral theorem, a Hermitian matrix of rank $\leq r$ can be written as $x=\sum_{s\leq r}\lambda_s\xi_s\xi_s^*$, with $\xi_1,\ldots,\xi_r$ orthonormal, and so lies over the point of the Grassmannian given by the span of these vectors, at least when its rank is exactly $r$. In the free setting, we are led in this way to the following definition:

\begin{definition}
The free Hermitian determinantal variety $D^+_{rN}$, with $1\leq r\leq N$, is given by
$$C(D^+_{rN})=C^*\left(x_{ij},p_{ij}\Big|x=x^*,\ \sum_{ij}x_{ij}x_{ij}^*=1,\ p=p^*=p^2,\ Tr(p)=r,\ px=x\right)$$
with $x,p\in M_N(A)$, the condition $x=x^*$ meaning $x_{ij}^*=x_{ji}$, for any $i,j$.
\end{definition}

Observe that by taking adjoints, the relation $px=x$ gives $xp=x$ as well, so that we have $x=pxp$, meaning that $x$ lives inside the projection $p$. Classically, we obtain the pairs $(x,\Lambda)$ with $x$ a Hermitian matrix of norm 1 supported by a complex subspace $\Lambda$ of dimension $r$, which form a sphere bundle over the complex Grassmannian $Gr^\mathbb C_{rN}$, of real dimension $2r(N-r)+r^2-1$, resolving the Hermitian matrices of rank $\leq r$ and norm 1. The program {\tt leoncavallo.py} checks the relations in Theorem \ref{thm:determinantal} (2,3) and in Theorem \ref{thm:hermitian} below, in random matrix models with entries in $M_K(\mathbb C)$, including a noncommutative model of $O_N^+$, given by $u_{ij}=o_{ij}g_j$, with $o\in O_N$, and $g_i$ random self-adjoint unitaries. We have:

\begin{theorem}
\label{thm:hermitian}
The free Hermitian determinantal varieties have the following properties:
\begin{enumerate}
\item We have an embedding $Gr^+_{rN}\subset D^+_{rN}$, given by $x=p/\sqrt{r}$.

\item At $r=1$, we have a map $S^{N-1}_{\mathbb R,+}\to D^+_{1N}$, given by $x_{ij}=p_{ij}=z_iz_j$.

\item The quantum group $O_N^+$ acts on $D^+_{rN}$, via $x\to u^txu$ and $p\to u^tpu$.
\end{enumerate}
Moreover, the maps in $(1)$ and $(2)$ are equivariant with respect to $O_N^+$.
\end{theorem}

\begin{proof}
(1) We have $\sum_{ij}p_{ij}p_{ij}^*=Tr(pp^*)=Tr(p)=r$, so $x=p/\sqrt{r}$ satisfies all the relations, with $px=p^2/\sqrt{r}=x$, and the corresponding morphism of algebras is surjective, because its image contains the coordinates $p_{ij}$.

(2) This follows from (1), by composing with the quotient map $S^{N-1}_{\mathbb R,+}\to P^{N-1}_+$.

(3) The coaction is $x_{ij}\to\sum_{kl}x_{kl}\otimes u_{ki}u_{lj}$, and similarly for $p$. The relations for $p$ are preserved as in section 4. Since the entries of $u$ are self-adjoint, so is $x$, and by using $uu^t=1$ we obtain $\sum_{ij}x_{ij}x_{ij}^*\to1\otimes1$, and $px=x$ is preserved, as in the rectangular case. The equivariance assertion is clear.
\end{proof}

\section{Stiefel and integration}

Following \cite{ba1}, consider the free partial isometries $u\in M_{M\times N}(A)$, with self-adjoint entries, such that $uu^t$ is a projection of trace $L$. At $M=L$, Lemma \ref{lem:trace} applied to $1-uu^t$ shows that this condition simply reads $uu^t=1_L$. Thus, we can formulate:

\begin{definition}
The free real and complex Stiefel manifolds are given by
\begin{eqnarray*}
C(St^+_{LN})&=&C^*\left((u_{ri})_{r\leq L,i\leq N}\Big|u=\bar{u},uu^t=1_L\right)\\
C(St^+_{LN,\mathbb C})&=&C^*\left((u_{ri})_{r\leq L,i\leq N}\Big|uu^*=\bar{u}u^t=1_L\right)
\end{eqnarray*}
that is, $St^+_{LN}=O_{LN}^{L+}$ and $St^+_{LN,\mathbb C}=U_{LN}^{L+}$, in the notations of \cite{ba1}.
\end{definition}

At $L=1$ we obtain the free spheres, and at $L=N$ we obtain $O_N^+$ and $U_N^+$, because Lemma \ref{lem:trace} applied to $1-u^tu$, respectively to $1-u^*u$ and $1-u^t\bar{u}$, gives the other unitarity relations. Classically, the Grassmannian is the quotient of the Stiefel manifold by $O_L$, via $U\to U^tU$, and in terms of dimensions, $St_{LN}=O_N/O_{N-L}$ has dimension $NL-L(L+1)/2$, while $Gr^\mathbb R_{LN}$ has dimension $L(N-L)$, the difference being $\dim O_L=L(L-1)/2$. In the free case, we have the maps $p_{ij}=\sum_ru_{ri}u_{rj}$ and $p_{ij}=\sum_ru_{ri}^*u_{rj}$, which are the $k=2$ case of Theorem \ref{thm:flagmaps}, and whose images are invariant under $O_L^+$, respectively $U_L^+$, acting via $u_{ri}\to\sum_sg_{rs}\otimes u_{si}$.

\bigskip

Let us discuss now the structure of the free Grassmannians, viewed inside $O_N^+$. We use the representation theory of $O_N^+$ from \cite{ba1}, at $N\geq2$, with the irreducible corepresentations $r_k\subset u^{\otimes k}$ being given by
$$r_k=\left\{\xi\in(\mathbb C^N)^{\otimes k}\Big|\textstyle\sum_x\xi_{\ldots xx\ldots}=0,\ \forall\ {\rm adjacent\ positions}\right\}$$
and with $P_k$ being the orthogonal projection onto $r_k$. For $\xi,\eta\in(\mathbb C^N)^{\otimes k}$ we use the coefficients $c_{\xi,\eta}=\sum_{IJ}\bar{\xi}_I\eta_J(u^{\otimes k})_{IJ}$, and similarly for $r_k$. Let $A_{LN}\subset C(O_N^+)$ be the $*$-algebra generated by the variables $p_{ij}=\sum_{r\leq L}u_{ri}u_{rj}$. With $\psi=\sum_{r\leq L}e_r\otimes e_r$, the products of $m$ such variables are the coefficients $c_{\psi^{\otimes m},\eta}$. We have:

\begin{lemma}
\label{lem:paths}
For $T\in Hom(r_j,u^{\otimes 2m})$ we have $T^*\psi^{\otimes m}\in\mathbb C\cdot P_j\psi^{\otimes j/2}$, with $j$ even.
\end{lemma}

\begin{proof}
By easiness, $T$ is a linear combination of maps $T_\pi P_j$, with $\pi$ being noncrossing pairings having $j$ through strings and caps on top, the caps not enclosing through strings. The vector $T_\pi^*\psi^{\otimes m}$ is computed by contracting along the caps, and the graph formed by the caps and by the edges $(2l-1,2l)$ of $\psi^{\otimes m}$ consists of cycles, contributing factors $L$, and of paths between through points. Since the caps stay in the gaps between consecutive through points, a path starting at a through point can never cross another through point, so the paths connect the through points in the order $(1,2),(3,4),\ldots$, and we obtain a multiple of $\psi^{\otimes j/2}$.
\end{proof}

This leads to the following result, which includes the first part of Theorem C:

\begin{theorem}
\label{thm:spectrum}
At $1\leq L\leq N-1$, the algebra $A_{LN}$ decomposes as
$$A_{LN}=\bigoplus_{k\geq0}C_k\quad,\quad C_k=\left\{c^{r_{2k}}_{\xi_k,\eta}\Big|\eta\in r_{2k}\right\}$$
with $\xi_k=P_{2k}\psi^{\otimes k}\neq0$. Thus $Gr^+_{LN}$ is multiplicity free, with spectrum $r_0,r_2,r_4,\ldots$
\end{theorem}

\begin{proof}
By Peter-Weyl, the component of $c_{\psi^{\otimes m},\eta}$ on $r_j$ is a sum of coefficients $c^{r_j}_{T^*\psi^{\otimes m},T^*\eta}$, with $T\in Hom(r_j,u^{\otimes2m})$, so by Lemma \ref{lem:paths} we have $A_{LN}\subset\bigoplus_kC_k$. Conversely, $c_{\psi^{\otimes k},P_{2k}\eta}=c^{r_{2k}}_{\xi_k,\eta}$, so $C_k\subset A_{LN}$. Finally, with $g=e_1+ie_N$ we have $g^{\otimes2k}\in r_{2k}$, because $\sum_xg_x^2=0$, and
$$<\psi^{\otimes k},g^{\otimes 2k}>=\left(\sum_{r\leq L}\bar{g}_r^2\right)^k=1$$
because $L<N$. Thus $\xi_k\neq0$, and this finishes the proof.
\end{proof}

For instance at $N=3$ we have $\dim r_{2k}=U_{2k}(3)=F_{4k+2}$, the Fibonacci numbers, so the dimensions of the summands are $1,8,55,377,\ldots$, growing like $\gamma^{2k}$, with $\gamma=(3+\sqrt{5})/2$. Classically, $Gr_{LN}$ is a symmetric space of rank $\min(L,N-L)$, which is multiplicity free, with spectrum indexed by the partitions having at most $\min(L,N-L)$ rows, see \cite{hel}. Thus, the free Grassmannians behave like rank one spaces. Also, since $\dim r_k$ grows like $\rho^k$, with $\rho+\rho^{-1}=N$, at $N\geq3$ these spaces are infinite dimensional in the Weyl law sense, as the free spheres. The space $C_1$ is spanned by the variables $p_{ij}-\frac{L}{N}\delta_{ij}$, because $\psi-P_2\psi=\frac{L}{N}\sum_xe_x\otimes e_x$, and this justifies the eigenspace claim in Theorem \ref{thm:plucker}, the Brownian eigenvalues $\lambda_{2k}^{group}$ from section 8 being increasing in $k$. As for the Laplacians, any operator on $A_{LN}$ commuting with the coaction of $O_N^+$ acts by scalars on the spaces $C_k$, and these Laplacians, together with the heat equation and the free analogue of the Weyl law, will be discussed in section 8.

\bigskip

Regarding now the quotient space $O_N^+/(O_L^+\,\widehat{*}\,O_{N-L}^+)$, this contains $A_{LN}$, as well as the following nested elements, which are invariant, by using $g^tg=1$ twice:
$$X_{abcd}=\sum_{r,s\leq L}u_{ra}u_{sb}u_{sc}u_{rd}$$

Classically $X_{abcd}=p_{ad}p_{bc}$. In the free case, we have the following result, which is the second part of Theorem C, in agreement with the situation for the free spheres \cite{ba1}, \cite{bss}:

\begin{theorem}
\label{thm:notgh}
At $N\geq3$ and $1\leq L\leq N-1$, the orbit algebra of $Gr^+_{LN}$, namely the image of $C(Gr^+_{LN})$ in $C(O_N^+)$, is a proper subalgebra of the algebra of coinvariants $C(O_N^+/(O_L^+\,\widehat{*}\,O_{N-L}^+))$. More generally, this orbit algebra is not of the form $C(O_N^+/H)$, for any closed subgroup $H\subset O_N^+$.
\end{theorem}

\begin{proof}
Assume first $2\leq L\leq N-1$. With $\psi'=\sum_{r,s\leq L}e_r\otimes e_s\otimes e_s\otimes e_r$, we have $X_{abcd}=c_{\psi',e_{abcd}}$, whose component on $r_4$ is $c^{r_4}_{P_4\psi',P_4e_{abcd}}$. If $X_{abcd}\in A_{LN}$ for any $a,b,c,d$, Theorem \ref{thm:spectrum} and Peter-Weyl would give $P_4\psi'\in\mathbb C\cdot P_4\psi^{\otimes2}$. But this is wrong, as shown by the following vector:
$$\zeta=e_1\otimes e_2\otimes e_2\otimes e_1-e_1\otimes e_N\otimes e_N\otimes e_1$$

Indeed, all the adjacent contractions of $\zeta$ vanish, so $\zeta\in r_4$, and we have $<\psi^{\otimes 2},\zeta>=0$ and $<\psi',\zeta>=1$. Since the isotypic projections are bounded, the same holds at the $C^*$-algebra level. At $L=1$, we use $p\to1-p$ and permute the rows, as to reduce the problem to the case $L=N-1\geq2$. For the last assertion, assume that our image is $O_N^+/H$. Then the multiplicity of $r_j$ is $\dim Fix_H(r_j)$, so by Theorem \ref{thm:spectrum} and by $u^{\otimes2}=r_0+r_2$ and $u^{\otimes4}=2r_0+3r_2+r_4$, we obtain $\dim Fix_H(u^{\otimes2})=2$ and $\dim Fix_H(u^{\otimes4})=6$, and so $\dim End_H(u)=2$ and $\dim End_H(u^{\otimes2})=6$. The first equality gives $u=a+b$, with $a,b$ irreducible and non-equivalent, with $\dim a\geq2$, because $N\geq3$. Thus $a\otimes a$, $a\otimes b$, $b\otimes a$ have dimension $\geq2$, and since each of them contains at most one copy of the trivial representation, each contains a nontrivial constituent. Together with the 2 copies of the trivial representation in $u^{\otimes2}$, this gives $\dim End_H(u^{\otimes2})\geq2^2+3=7$, a contradiction.
\end{proof}

As a numerical illustration, the program {\tt puccini.py} computes Gram matrices with respect to the Haar state, at $N=3$ and $L=2$, and finds that adding the elements $X_{abcd}$ to the span of the products of at most 3 entries of $p$ increases its dimension by $55=\dim r_4$, in agreement with the above proof.

\bigskip

The same method shows that, at $N\geq3$, the orbit algebra of the full flag manifold $Fl^+_N$ is a proper subalgebra of the algebra of $T_N^+$-coinvariants $C(O_N^+/T_N^+)$, from Theorem \ref{thm:flagmaps}. Indeed, the products of generators $q^s_{ij}=u_{si}u_{sj}$ have their row labels repeated in adjacent pairs, and the argument in the proof of Lemma \ref{lem:paths}, with fixed row labels, shows that their components on $r_4$ have left vectors in the span of the vectors $P_4e_{aabb}$. On the other hand, the coefficients $u_{1a}u_{2b}u_{2c}u_{1d}$, with left vector $e_{1221}$, are $T_N^+$-invariant, and $w=e_{1221}-e_{1331}$ belongs to $r_4$, is orthogonal to all the vectors $e_{aabb}$, and satisfies $<e_{1221},w>=1$. As a check, {\tt puccini.py} finds at $N=3$ that adding these coefficients to the span of the products of 2 generators of $Fl^+_3$ increases its dimension by $55=\dim r_4$.

\bigskip

Let us discuss now integration. By restricting the Haar functional of $O_N^+$, and using the Weingarten formula from \cite{ba1}, we obtain, with $\tau=\sqcap\sqcap\ldots\sqcap\in NC_2(2k)$:
$$\int_{Gr^+_{LN}}p_{i_1j_1}\ldots p_{i_kj_k}=\sum_{\pi,\sigma\in NC_2(2k)}L^{|\pi\vee\tau|}\delta_\sigma(i_1j_1\ldots i_kj_k)W_{2k,N}(\pi,\sigma)$$

Indeed, $p_{i_1j_1}\ldots p_{i_kj_k}$ is the sum of the products $u_{r_1i_1}u_{r_1j_1}\ldots u_{r_ki_k}u_{r_kj_k}$, and the sum over $r\in\{1,\ldots,L\}^k$ of $\delta_\pi(r_1r_1\ldots r_kr_k)$ is $L^{|\pi\vee\tau|}$. In particular, at $k=1,2$ we obtain:
$$\int_{Gr^+_{LN}}p_{11}=\frac{L}{N}\qquad,\qquad
\int_{Gr^+_{LN}}p_{11}^2=\frac{L(L+1)}{N(N+1)}$$

Classically, $p_{11}$ follows the beta law $\nu_{L/2,(N-L)/2}$ in the real case, and $\nu_{L,N-L}$ in the complex case, as explained in the introduction, with second moments $L(L+2)/N(N+2)$ and $L(L+1)/N(N+1)$, so as for the free spheres, the second moment coincides with the complex one. At $k\geq3$ the free moments are no longer given by product formulae, and for instance at $N=3$, $L=1$ the third moment is $2/21$, while the classical complex one is $1/10$. The program {\tt rossini.py} computes these moments, as exact rational numbers, by using the Weingarten functions from {\tt donizetti.py}. In the $N\to\infty$ limit, we have:

\begin{theorem}
\label{thm:poisson}
With $N\to\infty$ and $L$ fixed, the variable $Np_{11}\in C(Gr^+_{LN})$ converges in moments to the free Poisson law $\pi_L$, whose moments are $\sum_{\sigma\in NC(k)}L^{|\sigma|}$. In the classical real and complex cases, the limiting laws are respectively $\chi^2_{L1}=\mu_{L/2,2}$ and $\mu_{L1}$.
\end{theorem}

\begin{proof}
We have $W_{2k,N}(\pi,\sigma)=N^{-k}(\delta_{\pi\sigma}+O(N^{-1}))$, and so the moments of $Np_{11}$ converge to $\sum_{\pi\in NC_2(2k)}L^{|\pi\vee\tau|}$, which by fattening equals $\sum_{\sigma\in NC(k)}L^{|\sigma|}$, see \cite{nsp}. In the classical case, the same computation with $P_2(2k)$ and $\mathcal P_2(2k)$ gives the moments of $\chi^2_{L1}$ and of $\mu_{L1}$.
\end{proof}

Thus $\pi_L$, which is the law of a sum of $L$ squares of free semicircular variables, plays here the role of $\chi^2_{L1}$. For instance at $L=1$ the moments of $\pi_1$ are the Catalan numbers $1,2,5,14,\ldots$, while those of $\chi^2_{1,1}$ are $1,3,15,105,\ldots$ The program {\tt rossini.py} also checks this convergence numerically. The exact law of $p_{11}$ at fixed $N$ will be discussed in section 8. Observe also that in the regime $L=tN$, with $t\in(0,1)$ fixed, the above formulae give mean $t$ and variance $t(1-t)/(N+1)$, so $p_{11}$ converges to the constant $t$, with fluctuations of order $1/\sqrt{N}$.

\bigskip

Finally, regarding the flag manifolds, the diagonal entries $m_{si}=q^s_{ii}$ are of particular interest, and we have the following result, which will be generalized in section 10:

\begin{theorem}
Over the full flag manifold $Fl^+_N$, each diagonal entry $m_{si}$ has the same law as $p_{11}$ over $P^{N-1}_+$, with $Nm_{si}$ becoming free Poisson $\pi_1$, with $N\to\infty$. Moreover, any fixed number of entries in a row of $Nm$ are asymptotically free.
\end{theorem}

\begin{proof}
In the image of $O_N^+$ we have $m_{si}=u_{si}^2$, and the first assertion follows from the invariance of the Haar functional under the permutations of rows and columns, and from Theorem \ref{thm:poisson}. As for the asymptotic freeness, this follows from the fact that the rescaled coordinates of the free sphere become free semicircular with $N\to\infty$, see \cite{ba1}, \cite{bgo}, applied to the rows of $O_N^+$.
\end{proof}

\section{Spectral geometry}

In this section we discuss the heat equation and the Laplacians on the free Grassmannians, in relation with the work of Das, Franz and Wang \cite{dfw} on invariant Markov semigroups on quantum homogeneous spaces, and with the analysis of the free spheres from \cite{mfs}. We use the notations from section 7, and we denote by $\overline{A}_{LN}$ the norm closure of $A_{LN}$ in the reduced algebra $C_{red}(O_N^+)$, on which the Haar state is faithful.

\bigskip

As in \cite{mfs}, we use positive Laplacians $\mathcal L\geq0$, so that the heat semigroup is $T_s=e^{-s\mathcal L}$, with time $s\geq0$, the Markov generator is $-\mathcal L$, and the heat equation is $\partial_sf=-\mathcal Lf$. Recall that a convolution semigroup of states $(\varphi_s)$ on $O_N^+$ is called central when it acts on the coefficients of each irreducible corepresentation $r_j$ by a scalar, $\varphi_s(r_j)=e^{-s\lambda_j}1$. The basic example, coming from Cipriani, Franz and Kula \cite{cfk}, is the Brownian one, which normalized by $\lambda_1=1$ has eigenvalues given by the Chebyshev polynomials $U_j$ of the second kind, with $U_j(N)=\dim r_j$, as follows, see also \cite{mfs}, \cite{dfw}:
$$\lambda_j^{group}=\frac{NU_j'(N)}{U_j(N)}$$

The point is that the free Grassmannians are invariant under such semigroups, and that the resulting heat equation does not depend on $L$, as follows:

\begin{proposition}
\label{prop:central}
Any central convolution semigroup of states $(\varphi_s)$ on $O_N^+$, with eigenvalues $\lambda_j$, restricts into a symmetric Markov semigroup $T_s=(id\otimes\varphi_s)\Delta$ on $\overline{A}_{LN}$, preserving the Haar state, which acts on the spaces $C_k$ from Theorem \ref{thm:spectrum} by the scalars $e^{-s\lambda_{2k}}$. Thus, we have $T_s=e^{-s\mathcal L}$, with the positive Laplacian $\mathcal L$ on $Gr^+_{LN}$ acting on $C_k$, of dimension $U_{2k}(N)$, by $\lambda_{2k}$, independently of $1\leq L\leq N-1$.
\end{proposition}

\begin{proof}
We have $\Delta(A_{LN})\subset A_{LN}\otimes C(O_N^+)$, because $\Delta(p_{ij})=\sum_{ab}p_{ab}\otimes u_{ai}u_{bj}$. On a coefficient $c^{r_j}_{\xi,\eta}$ we have $(id\otimes\varphi_s)\Delta(c^{r_j}_{\xi,\eta})=e^{-s\lambda_j}c^{r_j}_{\xi,\eta}$, and Theorem \ref{thm:spectrum} gives the action on the spaces $C_k$. The positivity and the passage to the reduced algebra are as in \cite{mfs}, \cite{dfw}, by using the invariance of the Haar state.
\end{proof}

For instance at $N=3$ the first eigenvalues are $\lambda_2^{group}=9/4$ and $\lambda_4^{group}=54/11$, coming from $U_2(3)=8$, $U_2'(3)=6$, $U_4(3)=55$, $U_4'(3)=90$.

\bigskip

At $L=1$ the algebra $A_{1N}$ is the even part of the free sphere $S^{N-1}_{\mathbb R,+}$, generated by the variables $p_{ij}=z_iz_j$, and the spaces $C_k$ are the spaces $V_{2k}^\circ$ of spherical harmonics of even degree, from \cite{mfs}. Here there is a second natural heat semigroup, namely the intrinsic spherical semigroup of \cite{dfw}, whose eigenvalues $\lambda_r^+$ were computed in closed form in \cite{mfs}, and which preserves the even part. The 2 Laplacians differ already in degree 2 of the sphere, with $\lambda_2^{group}=2N^2/(N^2-1)$ and $\lambda_2^+=2N/(N-1)$. An intrinsic Laplacian on $Gr^+_{LN}$ with $2\leq L\leq N-2$, generalizing the intrinsic spherical Laplacian on $P^{N-1}_+$, would be very interesting.

\bigskip

In the classical case the heat trace $Tr(e^{-s\mathcal L})$ of a compact Riemannian manifold of dimension $d$, with $\mathcal L=-div\,grad$, is finite for any $s>0$, and behaves like $s^{-d/2}$, which is the Weyl law. In the free case the multiplicities grow exponentially, and the eigenvalues linearly, and we obtain instead a threshold, as follows:

\begin{theorem}
\label{thm:heattrace}
For $N\geq3$ and $1\leq L\leq N-1$, the heat trace of the Brownian semigroup on the free Grassmannian $Gr^+_{LN}$, namely $Z_N(s)=\sum_{k\geq0}U_{2k}(N)e^{-s\lambda_{2k}^{group}}$, is finite precisely for $s>s_N$, where $s_N=\frac{\sqrt{N^2-4}}{N}\log\gamma$, with $\gamma+\gamma^{-1}=N$, $\gamma>1$.
\end{theorem}

\begin{proof}
We have $U_j(N)=(\gamma^{j+1}-\gamma^{-j-1})/(\gamma-\gamma^{-1})$, and by logarithmic differentiation, using $d\gamma/dN=\gamma^2/(\gamma^2-1)$ and $\gamma-\gamma^{-1}=\sqrt{N^2-4}$, we obtain the following formula:
$$\lambda_j^{group}=\frac{N(j+1)}{\sqrt{N^2-4}}\left(1+\frac{2}{\gamma^{2j+2}-1}\right)-\frac{N^2}{N^2-4}$$

Thus $\lambda_j^{group}=N(j+1)/\sqrt{N^2-4}-N^2/(N^2-4)+O((j+1)\gamma^{-2j})$.

Thus, the terms of $Z_N(s)$ are equivalent to $\exp(2k(\log\gamma-sN/\sqrt{N^2-4}))$, up to a positive constant. These terms go to 0 exponentially fast when $s>s_N$, and tend to a positive constant at $s=s_N$, and this gives the result.
\end{proof}

Observe that with $N=2\cosh t$, the above formula reads $\lambda_j^{group}=\frac{N(j+1)\coth((j+1)t)}{2\sinh t}-\frac{N^2}{4\sinh^2t}$, which is increasing in $j$, because so is $x\coth x$, and at $N=2$ we have $U_j(2)=j+1$ and $U_j'(2)=j(j+1)(j+2)/6$, so $\lambda_j^{group}=j(j+2)/3$. For instance $s_3=0.71735\ldots$, $s_4=1.14052\ldots$, $s_5=1.43599\ldots$ and $s_{10}=2.24611\ldots$, and $s_N\sim\log N$ as $N\to\infty$. Physically, $Z_N(s)$ is a partition function at inverse temperature $s$, and the threshold $s_N$ corresponds to a maximal temperature $1/s_N$, as for the Hagedorn temperature \cite{hag}, see Fermi \cite{fth} for the thermodynamics background. For related asymptotic questions for the same semigroups, namely the cutoff phenomenon, see Freslon, Teyssier and Wang \cite{ftw}. The program {\tt bellini.py} computes the eigenvalues $\lambda_j^{group}$ and checks this threshold. A similar threshold, for the Hilbert-Schmidt norm of the intrinsic spherical heat operators, was found in \cite{mfs}.

\bigskip

Let us discuss now the law of the coordinates, refining the results from section 7. Recall from \cite{mfs} that on the free sphere $S^{N-1}_{\mathbb R,+}$ the squared radius of a block of $k$ coordinates is $E_k=\sum_{i\leq k}x_i^2$, and that at $N>2$ the law of a single coordinate $x_N$ has the following support, coming from the continuous $q$-ultraspherical polynomials, see \cite{bcz}, \cite{mfs}:
$$supp(law(x_N))=\left[-\frac{2}{\sqrt{N+2}},\frac{2}{\sqrt{N+2}}\right]$$

For $2\leq L\leq N-2$ the support of the law of $E_L$ is not known, and already the question of whether $0$ belongs to it, at $N=4$ and $L=2$, was raised in \cite{mfs}. In our setting, $0$ belonging to this support is a necessary condition for the coordinate characters of $Gr^+_{2,4}$ to survive in the reduced algebra. At $L=1$ and $L=N-1$, we have the following result, which shows that the reduced free projective spaces miss their coordinate points, as is the case for the free spheres and their poles, see \cite{mfs}:

\begin{theorem}
\label{thm:p11law}
The coordinate $p_{11}\in\overline{A}_{LN}$ has the same law as the block squared radius $E_L$ of the free sphere $S^{N-1}_{\mathbb R,+}$, with mean $L/N$ and variance $L(N-L)/N^2(N+1)$. Moreover, for $N\geq3$:
\begin{enumerate}
\item At $L=1$ the spectrum of $p_{11}$ is $[0,4/(N+2)]$.

\item At $L=N-1$ the spectrum of $p_{11}$ is $[(N-2)/(N+2),1]$.
\end{enumerate}
In particular, at $L=1$ and $L=N-1$, the coordinate characters, which are the torus fixed points of the classical Grassmannian, do not factor through $\overline{A}_{LN}$.
\end{theorem}

\begin{proof}
We have $p_{11}=\sum_{r\leq L}u_{r1}^2$. The antipode $S(u_{ij})=u_{ji}$ is antimultiplicative and leaves the Haar state invariant, and gives $S(p_{11}^k)=(\sum_{r\leq L}u_{1r}^2)^k$. Since the first row of $O_N^+$ is the free sphere, with its Haar state, the law of $p_{11}$ is that of $E_L$. The mean and variance follow from section 7, or from \cite{mfs}. At $L=1$ we have $p_{11}=x_1^2$, and at $L=N-1$ we have $p_{11}=1-x_N^2$, and the spectra follow from \cite{mfs}, the Haar state being faithful on the reduced algebra. Finally, a coordinate character maps $p_{ii}$ to 0 or 1, and these values are outside the spectra of the appropriate diagonal coordinates.
\end{proof}

\section{Arithmetic aspects}

In this section we discuss the free Grassmannians from the viewpoint of the free real algebraic geometry from \cite{fba}, and of the arithmetic aspects from \cite{oba}. Recall from \cite{fba} that a free real algebraic variety $X\subset\mathbb R^M_+$ is given by noncommutative polynomial relations between self-adjoint coordinates $x_1,\ldots,x_M$, and that $X$ is called archimedean when $R-\sum_ix_i^2$ is a sum of squares, modulo the relations, for some $R>0$. We have:

\begin{proposition}
\label{prop:archimedean}
The free Grassmannian $Gr^+_{LN}$ is an archimedean free real algebraic variety in $\mathbb R^{N^2}_+$, given by quadratic equations, in the self-adjoint coordinates
$$a_{ij}=\frac{p_{ij}+p_{ji}}{2}\ (i\leq j)\quad,\quad b_{ij}=\frac{p_{ij}-p_{ji}}{2i}\ (i<j)$$
and the same holds for the free flag manifolds, in $\mathbb R^{kN^2}_+$.
\end{proposition}

\begin{proof}
With $a$ symmetric and $b$ antisymmetric we have $p=a+ib$, and the condition $p=p^*$ is automatic. The relation $p^2=p$ is quadratic, and taking the operator real and imaginary parts $(f+f^*)/2$ and $(f-f^*)/2i$ of the entries $f=(p^2-p)_{ij}$ gives self-adjoint quadratic equations, while $Tr(p)=L$ is linear. Finally, from $Tr(p^2)=L$ we obtain, the cross terms cancelling by symmetry:
$$\sum_{i,j}(a_{ij}^2+b_{ij}^2)=\sum_{i,j}p_{ij}p_{ji}=L$$

Thus $L-\sum_{i\leq j}a_{ij}^2-\sum_{i<j}b_{ij}^2=\sum_{i<j}(a_{ij}^2+b_{ij}^2)$ modulo the relations, which is a sum of squares, and $Gr^+_{LN}$ is archimedean. For the flag manifolds the proof is similar.
\end{proof}

Regarding now the free quadrics in $P^{N-1}_+$ from section 3, given by $\sum_{ab}c_{ab}p_{ab}=0$, these are self-adjoint relations precisely when the matrix $H=(c_{ab})$ is Hermitian. On the free sphere they become the homogeneous free quadrics $\sum_{ab}H_{ab}z_az_b=0$, which are described in \cite{fba} by the decomposition $H=Q+iP$, with $Q$ real symmetric and $P$ real antisymmetric. Observe that in the projective setting the antisymmetric part $P$ is seen classically, on $P^{N-1}_\mathbb C$, where $p_{ab}=z_a\bar{z}_b$ and the quadric becomes the Hermitian quadric $\sum_{ab}H_{ab}z_a\bar{z}_b=0$, while on $P^{N-1}_\mathbb R$ only $Q$ is seen. For instance at $N=2$, with $Q=0$ and $P=(^{\ 0\ }_{-1}{\ }^1_0)$, the quadric is $p_{12}=p_{21}$. On $P^1_\mathbb C=S^2$ this is $Im(z_1\bar{z}_2)=0$, a great circle, namely $P^1_\mathbb R$, and the same holds on $P^1_+=P^1_\mathbb C$.

\bigskip

Let us discuss now the ground field. Following \cite{oba}, where the easy quantum groups were studied over arbitrary fields of characteristic 0, we can formulate:

\begin{definition}
Given a field $F$ of characteristic $0$, we let $\mathcal O(Gr^+_{LN})_F$ be the $F$-algebra generated by variables $p_{ij}$, with $i,j=1,\ldots,N$, with the relations $p^2=p$ and $Tr(p)=L$.
\end{definition}

Classically, the commutative version of this algebra is the coordinate algebra of the affine variety of idempotent matrices of trace $L$, which over $\mathbb C$ is the homogeneous space $GL_N/(GL_L\times GL_{N-L})$, an affine bundle over the Grassmannian. The free Grassmannian appears as a real form of this algebra, in the sense of \cite{oba}, as follows:

\begin{proposition}
\label{prop:realform}
The map $\iota(p_{ij})=p_{ji}$ extends into an involutive anti-automorphism of the algebra $\mathcal O(Gr^+_{LN})_\mathbb Q$. Over $\mathbb C$, the antilinear map $a\otimes z\to\iota(a)\otimes\bar{z}$ is an involution, making $p$ self-adjoint, so this is a real form in the sense of $*$-structures, and the enveloping $C^*$-algebra is $C(Gr^+_{LN})$. The same holds for the free flag manifolds.
\end{proposition}

\begin{proof}
The map $\iota$ is well-defined, because it maps the relation $(p^2-p)_{ij}=0$ into $(p^2-p)_{ji}=0$, and $Tr(p)=L$ into itself. The other assertions are clear from definitions, the enveloping $C^*$-algebra existing because the entries of a projection have norm $\leq1$, as in \cite{oba}.
\end{proof}

At $L=1$, $N=2$ the same computation as in section 1, which does not use the involution, shows that $\mathcal O(Gr^+_{1,2})_F$ is commutative, being the coordinate algebra of the affine quadric surface $ce=a(1-a)$, where $p=(^a_e{\ }^c_{1-a})$. At $L=1$ and $N\geq3$ these algebras are noncommutative, over any field $F$, as shown by the model $p_{ij}=g_ig_j/N$ from section 1, with $g_1,\ldots,g_N$ being the generators of $\mathbb Z_2^{*N}$.

Since the Weingarten matrices have rational entries, the Haar integrals over the free Grassmannians of the polynomials with rational coefficients in the projective coordinates are rational, for any $N\geq2$. More precisely, by the Weingarten formula from section 7, and by \cite{oba}, the integral of a product of $k$ projective coordinates belongs to $\mathbb Z[1/\det G_{2k,N}]$, where $G_{2k,N}$ is the Gram matrix of $O_N^+$, as in the formula $\int p_{11}^2=L(L+1)/N(N+1)$, whose denominator divides $\det G_{4,N}=N^2(N^2-1)$. In the $p$-adic setting, it was shown in \cite{oba} that the free Poisson law $\pi_1$ is not a bounded measure on $\mathbf Z_p$. The program {\tt pergolesi.py} computes the binomial moments of $\pi_t$ up to order 80, and finds that their $p$-adic valuations decrease roughly linearly for $t=1,2,3$, suggesting that the laws $\pi_L$ from Theorem \ref{thm:poisson} are not bounded either, while for $t=p^a$ with $a$ large, such as $t=81$ at $p=3$, these valuations stay $\geq0$ in this range. However, at higher order the same program finds $v_3(b_{138})=-1$ at $t=81$, and $v_5(b_{180})=-1$ at $t=125$, where $b_n$ are the binomial moments. It would be interesting to know if $\pi_t$ is a bounded measure for $v_p(t)$ large, as is the case for the Poisson law, by \cite{oba}.

\bigskip

Summarizing, we have the following result, which shows that the free Grassmannians fit with the free real algebraic geometry from \cite{fba}, and with the arithmetic from \cite{oba}:

\begin{theorem}
The free Grassmannians $Gr^+_{LN}$ are compact, archimedean intersections of free quadrics in $\mathbb R^{N^2}_+$, contained in a free sphere of radius $\sqrt{L}$, with coordinates $a_{ii}$ and $\sqrt{2}a_{ij},\sqrt{2}b_{ij}$ for $i<j$, they are real forms of the free affine varieties of idempotents of trace $L$, which are defined over $\mathbb Q$, and their Haar integrals of rational polynomials are rational.
\end{theorem}

\begin{proof}
This follows from Proposition \ref{prop:archimedean}, Proposition \ref{prop:realform}, and the above discussion.
\end{proof}

\section{Fermi, modernity}

This last section is dedicated to Enrico Fermi, whose work is at the origin of the Standard Model, and whose name appears, via the fermions, in the physical trace on its finite algebra. Besides the Fermi-Dirac statistics, Fermi developed in 1933-1934 the theory of the beta decay \cite{fer}, with a contact interaction between 4 fermions, namely the neutron, the proton, the electron and the neutrino, the latter being named by him. The decay rates were computed there via the perturbation theory formula going back to Dirac \cite{dir}, which gives the probability per unit time of a transition $i\to f$, induced by a perturbation $H$, with $\rho$ being the density of final states, as follows:
$$\Gamma_{i\to f}=\frac{2\pi}{\hbar}\,|<f|H|i>|^2\rho(E_f)$$

This formula was called a golden rule by Fermi, in his lecture notes \cite{fnp}, and is now known as Fermi's golden rule. It is a key ingredient of particle physics, the decay rates and the cross sections being all computed in this way, from the transition amplitudes.

\bigskip

As for Fermi's theory itself, after the work of Feynman and Gell-Mann \cite{fge} on its $V-A$ structure, it led to the electroweak theory of Glashow, Weinberg and Salam \cite{gla}, \cite{wei}, where the Fermi constant appears as $G_F/\sqrt{2}=g^2/8M_W^2$, in the low energy limit of the exchange of a $W$ boson, its measured value being $G_F\simeq1.166\cdot10^{-5}\,GeV^{-2}$. The mixing between the quark generations was introduced by Cabibbo \cite{cab}, and extended to 3 generations by Kobayashi and Maskawa \cite{kma}, while the neutrino mixing goes back to Pontecorvo \cite{pon}, a member of Fermi's group in Rome, and to Maki, Nakagawa and Sakata \cite{mns}. Thus, Italian physics is present at the origins of the Standard Model, and our purpose here is to show that the Italian projective geometry, made free, is present too.

\bigskip

We use the noncommutative geometry formulation of the Standard Model, due to Chamseddine and Connes \cite{ccs}, \cite{ccm}, and the quantum permutation group viewpoint from \cite{bop}. The finite part of the Standard Model has algebra $A_F=\mathbb C\oplus\mathbb H\oplus M_3(\mathbb C)$, and the complex algebra generated by it is $B_F=\mathbb C\oplus\mathbb C\oplus M_2(\mathbb C)\oplus M_3(\mathbb C)$, see \cite{bop}. The starting observation is that the pure states of $M_N(\mathbb C)$ are the density matrices $\rho$ which are rank one projections, $\rho=\rho^*=\rho^2$ and $Tr(\rho)=1$, so they form the complex projective space $P^{N-1}_\mathbb C$, see Bengtsson and Zyczkowski \cite{bzy}. Imposing the same relations on noncommuting coordinates gives $P^{N-1}_+$, which can be therefore regarded as the free space of pure states of $M_N(\mathbb C)$. We have:

\begin{theorem}
\label{thm:purestates}
For $N\geq2$, the quantum automorphism group of $M_N(\mathbb C)$ with its trace is the projective quantum isometry group $PO_N^+$ of the free projective space $P^{N-1}_+$, the free space of pure states of $M_N(\mathbb C)$, while the classical automorphism group $PU_N$ is the projective quantum isometry group of the space of pure states $P^{N-1}_\mathbb C$.
\end{theorem}

\begin{proof}
The quantum automorphism group of $M_N(\mathbb C)$, with its trace, is $PU_N^+$, see \cite{bop}, and we have $PU_N^+=PO_N^+$, see \cite{bfu}. Now by Theorem \ref{thm:projiso} this is the projective quantum isometry group of $P^{N-1}_+$, and the classical statement follows from Theorem \ref{thm:projiso} too.
\end{proof}

At $N=2$ we have $PO_2^+=PU_2^+=SO_3$, which is the quantum automorphism group of $M_2(\mathbb C)$, see \cite{ba1}, \cite{bop}, in agreement with the fact that the free projective line is classical, $P^1_+=P^1_\mathbb C=S^2$, as explained in section 1. At $N=3$ we obtain the color quantum group $PU_3^+$. Thus, the results in \cite{bop}, coming from the computation of Bhowmick, D'Andrea and Dabrowski \cite{bdd}, and from the twisting results from \cite{bop}, can be reformulated as follows:

\begin{theorem}
\label{thm:smpure}
With the trace coming from the fermions, the quantum automorphism group of $B_F$ is the free product of the projective quantum isometry groups of the free projective line $P^1_+$ and of the free projective plane $P^2_+$, with algebra $C(SO_3)*C(PU_3^+)$, and the color parts of the quantum isometry groups from \cite{bdd} are the projective quantum isometry groups of $P^2_+$ in the complex setting, and of $P^2_\mathbb C$ in the real setting.
\end{theorem}

\begin{proof}
The first assertion is from \cite{bop}, reformulated via Theorem \ref{thm:purestates}. The second one follows from the computation in \cite{bdd}, under the assumptions there on the mass matrices, which gives the color parts $PU_3^+$ in the complex setting and $PU_3$ in the real setting, and from Theorem \ref{thm:purestates}.
\end{proof}

Thus, the dichotomy between the complex and real settings in \cite{bdd} is the dichotomy between the free projective plane $P^2_+$ and the classical one $P^2_\mathbb C$. This is in agreement with the philosophy that the geometry should become free at very small scales. Observe also that at $N=2$ we have $P^1_+=P^1_\mathbb C=S^2$, which is the Bloch sphere of the pure states of a qubit, so the weak part sees no freeness at the level of pure states, while the color part sees the free projective plane $P^2_+$, which is the free space of pure states of a qutrit. Regarding the dynamics, the spectral action \cite{ccs} is built from the heat type traces, and Theorem \ref{thm:heattrace} shows that the trace $Tr(e^{-\mathcal L/\Lambda^2})$ on the free Grassmannians, with $\mathcal L$ being the Brownian Laplacian, is finite precisely below the energy scale $\Lambda_N=s_N^{-1/2}$. This is formally analogous to the Hagedorn temperature \cite{hag}, which appears in hadron physics and in string theory, as a consequence of the exponential growth of the number of states.

\bigskip

The other objects in this paper have interpretations of the same type. The density matrices with prescribed spectrum $\lambda_1,\ldots,\lambda_k\geq0$, with multiplicities $d_1,\ldots,d_k$ and with $\sum_sd_s\lambda_s=1$, form the unitary orbits which are the flag manifolds, see \cite{bzy}, and Proposition \ref{prop:isospectral} shows that the free flag manifolds $Fl^+_d$ are the free spaces of density matrices with prescribed spectrum. In particular the free Grassmannians correspond to the normalized projections $\rho=p/L$, which are the maximally mixed states on the subspaces of dimension $L$.

\bigskip

Regarding now the mixing, consider the mass eigenbases $(e_s)$ and $(f_t)$ of the up and down quarks, ordered by the masses. By Fermi's golden rule, the rates of the weak transitions between the quarks $s$ and $t$ are proportional to the transition probabilities $|<e_s,f_t>|^2=|V_{st}|^2$, where $V_{st}=<f_t,e_s>$ is the CKM matrix, and the same holds for the leptons, with the PMNS matrix, at the level of the charged-current vertices, the oscillation probabilities involving as well propagation phases, see \cite{bop}. Thus, the eigenbases are points of the full flag manifold $Fl_3$, and the unistochastic matrix $(|V_{st}|^2)$ measures their relative position. In terms of the projections $q^s,r^t$ onto $e_s,f_t$, we have $|<e_s,f_t>|^2=Tr(q^sr^t)=\sum_{ij}q^s_{ij}r^t_{ji}$, and in the real case this is $\sum_{ij}q^s_{ij}r^t_{ij}$. On complex points, $\sum_{ij}q^s_{ij}r^t_{ij}=|<e_s,\bar{f}_t>|^2$ is the overlap with the conjugate flag, and this latter quantity, which keeps the order of the indices, is the one having a good free behavior. We are led in this way to:

\begin{proposition}
\label{prop:transition}
Given free flags $q\in Fl^+_d$ and $r\in Fl^+_e$, with $d,e$ being compositions of $N$, the entries of the free transition matrix
$$\theta_{st}=\sum_{ij}q^s_{ij}\otimes r^t_{ij}\in C(Fl^+_d\times Fl^+_e)$$
are positive, and we have $\sum_t\theta_{st}=d_s$ and $\sum_s\theta_{st}=e_t$. When $r$ is the coordinate full flag, $r^t=e_{tt}$, we recover the matrix of diagonal entries $m_{st}=q^s_{tt}$ from section $5$.
\end{proposition}

\begin{proof}
Consider the elements $x_{ab}=\sum_jq^s_{aj}\otimes r^t_{bj}$, whose adjoints are $x_{ab}^*=\sum_iq^s_{ia}\otimes r^t_{ib}$. Since the multiplication on $C(Fl^+_d)\otimes C(Fl^+_e)$ is done legwise, by using $(q^s)^2=q^s$ and $(r^t)^2=r^t$ we obtain:
$$\sum_{ab}x_{ab}^*x_{ab}=\sum_{ij}\left(\sum_aq^s_{ia}q^s_{aj}\right)\otimes\left(\sum_br^t_{ib}r^t_{bj}\right)=\theta_{st}$$

Thus $\theta_{st}\geq0$. Also, $\sum_tr^t=1$ gives $\sum_t\theta_{st}=Tr(q^s)=d_s$, and similarly for the columns. The last assertion is clear from definitions.
\end{proof}

Classically, in the real case, $\theta$ is the orthostochastic matrix $(<e_s,f_t>^2)$, and in the complex case, $\theta$ is the unistochastic matrix $(|<e_s,\bar{f}_t>|^2)$. Observe that Theorem \ref{thm:realcomplex} identifies the laws of single flags, but not these pointwise observables, and for instance the free third moments of $\sum_{ij}q_{ij}r_{ij}$ and $\sum_{ij}q_{ij}r_{ji}$, at $N=3$ and rank one, are $2/21$ and $9/98$. The relative position of 2 complex flags is however not determined by such a matrix, the phases of $V$ modulo the torus carrying the CP violation, which is measured by the Jarlskog invariant $J=Im(V_{11}V_{22}\bar{V}_{12}\bar{V}_{21})$, see \cite{jar}. Observe that we have, with the above notations:
$$Tr(q^1r^1q^2r^2)=V_{11}\bar{V}_{21}V_{22}\bar{V}_{12}$$

In the free case, where we have commuting legs, this suggests considering the following variables, the imaginary part $J$ being a free analogue of the Jarlskog invariant:
$$X=\sum_{ijkl}q^1_{ij}q^2_{kl}\otimes r^1_{jk}r^2_{li}\quad,\quad J=\frac{X-X^*}{2i}$$


Regarding now the integration, we use on $Fl^+_d$ the integration coming from $O_N^+$, via Theorem \ref{thm:flagmaps}, as in section 7, and the product integration on $Fl^+_d\times Fl^+_e$, corresponding to 2 independent random flags. Classically, the relative position of 2 independent random real flags is given by a random orthogonal matrix. In the free case, the classical coset description fails, by section 7, but we have the following result, which generalizes the last theorem in section 7, and which is checked by brute force at small $N$ by the program {\tt monteverdi.py}:

\begin{theorem}
\label{thm:transitionlaw}
The family $(\theta_{st})$ has the same joint law as the family of block variables
$$\eta_{st}=\sum_{a\in B_s,b\in C_t}u_{ab}^2\in C(O_N^+)$$
with $B_s,C_t$ being the blocks of $d,e$. In particular, for $2$ independent free full flags, $\theta$ has the same law as the matrix $(u_{st}^2)$, and for $2$ independent points of $P^{N-1}_+$, the transition probability $\theta_{11}$ has the same law as the coordinate $p_{11}$.
\end{theorem}

\begin{proof}
By the Weingarten formula, with $\alpha=a_1a_1\ldots a_ka_k$, $\beta=b_1b_1\ldots b_kb_k$ and $\gamma=i_1j_1\ldots i_kj_k$, and with the sums being over $a_l\in B_{s_l}$, $b_l\in C_{t_l}$ and over the indices $i_l,j_l$, we have:
$$\int\theta_{s_1t_1}\ldots\theta_{s_kt_k}=\sum_{\alpha\beta\gamma}\sum_{\pi\sigma\pi'\sigma'}\delta_\pi(\alpha)\delta_\sigma(\gamma)W(\pi,\sigma)\delta_{\pi'}(\beta)\delta_{\sigma'}(\gamma)W(\pi',\sigma')$$

The sum over $\gamma$ of $\delta_\sigma(\gamma)\delta_{\sigma'}(\gamma)$ is $G(\sigma,\sigma')$, and since $W=G^{-1}$ is symmetric, we have $WGW=W$. Thus our integral is $\sum_{\alpha\beta}\sum_{\pi\pi'}\delta_\pi(\alpha)\delta_{\pi'}(\beta)W(\pi,\pi')$, which by the Weingarten formula equals $\int\eta_{s_1t_1}\ldots\eta_{s_kt_k}$. The particular cases follow from this, by using the fact that $p_{11}=u_{11}^2$ over $P^{N-1}_+$.
\end{proof}

Alternatively, the map $\tau(u_{ij})=u_{ji}$ extends into a $*$-automorphism of $C(O_N^+)$, which reverses the comultiplication, and so preserves the Haar state. With $\beta=(id\otimes\tau)\Delta$ we have $\beta(u_{ab})=\sum_iu_{ai}\otimes u_{bi}$, so $\theta_{st}=\beta(\eta_{st})$ in the orbit models, and $(h\otimes h)\beta=h$ gives the equality of all joint $*$-moments. This explains as well the order of the indices in the definition of $\theta$.

Observe that, by Theorem \ref{thm:transitionlaw} and by the last theorem in section 7, the free transition probabilities $\theta_{st}$ between 2 independent free full flags become, with $N\to\infty$ and after rescaling by $N$, free Poisson of parameter 1.

\bigskip

Regarding now the free Jarlskog variable $J$, in the classical real case we have $J=0$, because $Tr(q^1r^1q^2r^2)$ is then real. In the free case, where the real and complex cases coincide, the situation is different, and a CP-odd component is present even in the real picture, as shown by the following result, whose moments are computed exactly by the program {\tt vivaldi.py}, which as a check finds as well $J=0$ in the classical real case. In the classical complex case, the law of $J$ for a Haar random unitary in $U_3$ was computed by Dunkl and Zyczkowski \cite{dzy}. We have:

\begin{theorem}
\label{thm:jarlskog}
For $N\geq2$, and $2$ independent free full flags $q,r\in Fl^+_N$, the free Jarlskog variable $J$ is centered, with second moment given by
$$\int J^2=\frac{(N-2)P(N)}{2N(N^2-1)(N^2-2)^2(N^4-3N^2+1)^2}$$
where $P(N)=N^{10}-9N^8-2N^7+32N^6+8N^5-46N^4-18N^3+29N^2+14N-4$, so that $J\neq0$, for any $N\geq3$.
\end{theorem}

\begin{proof}
Every integral of a word in the variables $u_{ij}$ is real, so the integral of $(X-X^*)^{2m+1}$ is both real and purely imaginary, and so all the odd moments of $J$ vanish. Also, {\tt vivaldi.py} finds $\int X=-1/N(N^2-1)$. As for $\int J^2=-\int(X-X^*)^2/4$, the moments of order 2 of $X,X^*$ are computed as in the proof of Theorem \ref{thm:transitionlaw}, with the difference that the indices of the second leg are cyclically shifted with respect to those of the first leg, so that $G(\sigma,\sigma')$ is replaced by $N^{c(\sigma,\sigma')}$, with $c(\sigma,\sigma')$ being the number of connected components of the graph on the 8 indices formed by $\sigma,\sigma'\in NC_2(8)$. The program {\tt vivaldi.py} computes these sums exactly, for $N=2,\ldots,130$, and finds the above formula. This proves the formula for any $N$. Indeed, $G_{8,N}$ is a $14\times14$ matrix with entries of degree $\leq4$, so its cofactors have degree $\leq52$, and the loop factors have degree $\leq8$, so the result is of the form $A(N)/\det(G_{8,N})^2$, with $\deg A\leq112$. Since the proposed numerator and denominator have degrees 11 and 15, clearing denominators gives a polynomial of degree $\leq127$, vanishing at the 129 points $N=2,\ldots,130$. Finally, $P(t+3)$ has positive coefficients, so $\int J^2>0$ for $N\geq3$, and so $J\neq0$. For instance at $N=3,4,5$ we obtain the values $3397/1422960$, $138209/128422140$, $6588541/12848394320$.
\end{proof}

\begin{thebibliography}{99}
\addcontentsline{toc}{section}{References}
\setlength{\itemsep}{1pt}
\setlength{\parsep}{0pt}
\raggedright

\bibitem{ags}F. Arici, F. Galuppi and T. Gateva-Ivanova, Veronese and Segre morphisms between non-commutative projective spaces, {\em Eur. J. Math.} {\bf 8} (2022), 235--273.

\bibitem{bfu}T. Banica, The free unitary compact quantum group, {\em Comm. Math. Phys.} {\bf 190} (1997), 143--172.

\bibitem{ba1}T. Banica, Affine noncommutative geometry, {\color{black}arxiv:2012.10973} (2026).

\bibitem{bcz}T. Banica, B. Collins and P. Zinn-Justin, Spectral analysis of the free orthogonal matrix, {\em Int. Math. Res. Not. IMRN} {\bf 17} (2009), 3286--3309.

\bibitem{bgo}T. Banica and D. Goswami, Quantum isometries and noncommutative spheres, {\em Comm. Math. Phys.} {\bf 298} (2010), 343--356.

\bibitem{bop}T. Banica and C. Opus, Quantum permutations and the Standard Model algebra, \href{https://hexagonmath.org/2611.00722v1}{\textcolor{black}{\mbox{hexagon:2611.00722}}} (2026).

\bibitem{bss}T. Banica, A. Skalski and P.M. So\l tan, Noncommutative homogeneous spaces: the matrix case, {\em J. Geom. Phys.} {\bf 62} (2012), 1451--1466.

\bibitem{bcf}I. Baraquin, G. C\'ebron, U. Franz, L. Maassen and M. Weber, De Finetti theorems for the unitary dual group, {\em SIGMA} {\bf 18} (2022), 067.

\bibitem{bch}Y. Bazlov and R. Chen, Quantum Segre maps via cocycle twists, {\color{black}arxiv:2501.08942}.

\bibitem{bzy}I. Bengtsson and K. Zyczkowski, Geometry of quantum states, Cambridge Univ. Press (2006).

\bibitem{bdd}J. Bhowmick, F. D'Andrea and L. Dabrowski, Quantum isometries of the finite noncommutative geometry of the Standard Model, {\em Comm. Math. Phys.} {\bf 307} (2011), 101--131.

\bibitem{bro}L.G. Brown, Ext of certain free product $C^*$-algebras, {\em J. Operator Theory} {\bf 6} (1981), 135--141.

\bibitem{cab}N. Cabibbo, Unitary symmetry and leptonic decays, {\em Phys. Rev. Lett.} {\bf 10} (1963), 531--533.

\bibitem{ccs}A.H. Chamseddine and A. Connes, The spectral action principle, {\em Comm. Math. Phys.} {\bf 186} (1997), 731--750.

\bibitem{ccm}A.H. Chamseddine, A. Connes and M. Marcolli, Gravity and the standard model with neutrino mixing, {\em Adv. Theor. Math. Phys.} {\bf 11} (2007), 991--1089.

\bibitem{cfk}F. Cipriani, U. Franz and A. Kula, Symmetries of L\'evy processes on compact quantum groups, their Markov semigroups and potential theory, {\em J. Funct. Anal.} {\bf 266} (2014), 2789--2844.

\bibitem{cls}L. Cirio, G. Landi and R.J. Szabo, Algebraic deformations of toric varieties II: noncommutative instantons, {\em Adv. Theor. Math. Phys.} {\bf 15} (2011), 1817--1907.

\bibitem{mfs}Codex and Teo, The magic of free spheres, \href{https://hexagonmath.org/2611.00876v1}{\textcolor{black}{\mbox{hexagon:2611.00876}}} (2026).

\bibitem{co2}A. Connes, A unitary invariant in Riemannian geometry, {\em Int. J. Geom. Methods Mod. Phys.} {\bf 5} (2008), 1215--1242.

\bibitem{ddl}F. D'Andrea, L. Dabrowski and G. Landi, The noncommutative geometry of the quantum projective plane, {\em Rev. Math. Phys.} {\bf 20} (2008), 979--1006.

\bibitem{dfw}B. Das, U. Franz and X. Wang, Invariant Markov semigroups on quantum homogeneous spaces, {\em J. Noncommut. Geom.} {\bf 15} (2021), 531--580.

\bibitem{dst}M.S. Dijkhuizen and J.V. Stokman, Quantized flag manifolds and irreducible $*$-representations, {\em Comm. Math. Phys.} {\bf 203} (1999), 297--324.

\bibitem{dir}P.A.M. Dirac, The quantum theory of the emission and absorption of radiation, {\em Proc. Roy. Soc. A} {\bf 114} (1927), 243--265.

\bibitem{dzy}C. Dunkl and K. Zyczkowski, Volume of the set of unistochastic matrices of order 3 and the mean Jarlskog invariant, {\em J. Math. Phys.} {\bf 50} (2009), 123521.

\bibitem{fba}C. Fable and T. Banica, Noncommutative curves and surfaces in small degree, \href{https://hexagonmath.org/2611.00818v1}{\textcolor{black}{\mbox{hexagon:2611.00818}}} (2026).

\bibitem{fer}E. Fermi, Versuch einer Theorie der $\beta$-Strahlen I, {\em Z. Phys.} {\bf 88} (1934), 161--177.

\bibitem{fnp}E. Fermi, Nuclear physics, Univ. of Chicago Press (1950).

\bibitem{fth}E. Fermi, Thermodynamics, Dover (1956).

\bibitem{fge}R.P. Feynman and M. Gell-Mann, Theory of the Fermi interaction, {\em Phys. Rev.} {\bf 109} (1958), 193--198.

\bibitem{fio}R. Fioresi, Quantization of the Grassmannian manifold, {\em J. Algebra} {\bf 214} (1999), 418--447.

\bibitem{ftw}A. Freslon, L. Teyssier and S. Wang, Cutoff profiles for quantum L\'evy processes and quantum random transpositions, {\em Probab. Theory Related Fields} {\bf 183} (2022), 1285--1327.

\bibitem{fha}W. Fulton and J. Harris, Representation theory: a first course, Springer (1991).

\bibitem{gla}S.L. Glashow, Partial-symmetries of weak interactions, {\em Nucl. Phys.} {\bf 22} (1961), 579--588.

\bibitem{go1}D. Goswami, Quantum group of isometries in classical and noncommutative geometry, {\em Comm. Math. Phys.} {\bf 285} (2009), 141--160.

\bibitem{gjo}D. Goswami and S. Joardar, Non-existence of faithful isometric action of compact quantum groups on compact, connected Riemannian manifolds, {\em Geom. Funct. Anal.} {\bf 28} (2018), 146--178.

\bibitem{hag}R. Hagedorn, Statistical thermodynamics of strong interactions at high energies, {\em Nuovo Cimento Suppl.} {\bf 3} (1965), 147--186.

\bibitem{har}J. Harris, Algebraic geometry, Springer (1992).

\bibitem{hel}S. Helgason, Groups and geometric analysis, Academic Press (1984).

\bibitem{hsz}J.H. Hong and W. Szymanski, Quantum spheres and projective spaces as graph algebras, {\em Comm. Math. Phys.} {\bf 232} (2002), 157--188.

\bibitem{ito}H. Ito, NC functions over the nc Grassmannian, {\color{black}arxiv:2601.18041}.

\bibitem{jar}C. Jarlskog, Commutator of the quark mass matrices in the standard electroweak model and a measure of maximal CP nonconservation, {\em Phys. Rev. Lett.} {\bf 55} (1985), 1039--1042.

\bibitem{kma}M. Kobayashi and T. Maskawa, CP-violation in the renormalizable theory of weak interaction, {\em Progr. Theoret. Phys.} {\bf 49} (1973), 652--657.

\bibitem{krs}S. Kruglyak, V. Rabanovich and Yu. Samoilenko, Decomposition of a scalar matrix into a sum of orthogonal projections, {\em Linear Algebra Appl.} {\bf 370} (2003), 217--225.

\bibitem{lau}A. Lauve, Quantum- and quasi-Pl\"ucker coordinates, {\em J. Algebra} {\bf 296} (2006), 440--461.

\bibitem{mns}Z. Maki, M. Nakagawa and S. Sakata, Remarks on the unified model of elementary particles, {\em Progr. Theoret. Phys.} {\bf 28} (1962), 870--880.

\bibitem{mcc}K. McClanahan, $C^*$-algebras generated by elements of a unitary matrix, {\em J. Funct. Anal.} {\bf 107} (1992), 439--457.

\bibitem{nsp}A. Nica and R. Speicher, Lectures on the combinatorics of free probability, Cambridge Univ. Press (2006).

\bibitem{oba}C. Opus and T. Banica, Easiness over characteristic zero fields, \href{https://hexagonmath.org/2611.00831v1}{\textcolor{black}{\mbox{hexagon:2611.00831}}} (2026).

\bibitem{pon}B. Pontecorvo, Inverse beta processes and nonconservation of lepton charge, {\em Soviet Physics JETP} {\bf 7} (1958), 172--173.

\bibitem{sha}I.R. Shafarevich, Basic algebraic geometry, Springer (1974).

\bibitem{shu}T. Shulman, On universal $C^*$-algebras generated by $n$ projections with scalar sum, {\color{black}arxiv:0707.3053}.

\bibitem{tak}T. Takahashi, Minimal immersions of Riemannian manifolds, {\em J. Math. Soc. Japan} {\bf 18} (1966), 380--385.

\bibitem{voi}D.V. Voiculescu, Free analysis questions II: the Grassmannian completion and the series expansions at the origin, {\em J. Reine Angew. Math.} {\bf 645} (2010), 155--236.

\bibitem{wa1}S. Wang, Free products of compact quantum groups, {\em Comm. Math. Phys.} {\bf 167} (1995), 671--692.

\bibitem{wa2}S. Wang, Quantum symmetry groups of finite spaces, {\em Comm. Math. Phys.} {\bf 195} (1998), 195--211.

\bibitem{wei}S. Weinberg, A model of leptons, {\em Phys. Rev. Lett.} {\bf 19} (1967), 1264--1266.

\bibitem{wo1}S.L. Woronowicz, Compact matrix pseudogroups, {\em Comm. Math. Phys.} {\bf 111} (1987), 613--665.

\end{thebibliography}

\end{document}
